%% ELAPSE: Entropy-Linked Autonomous Propagation and Self-Erasure
%% TARGET: Physica A: Statistical Mechanics and its Applications — Q1, IF ≈ 3.3
%% SCOPE EMPHASIS for Physica A:
%%   - Statistical physics framing: phase transitions, criticality, scaling laws
%%   - Percolation remark elevated to a dedicated subsection "§4.4 Critical
%%     Threshold and Phase Transition near H_c"
%%   - Adversarial section REMOVED from main body (moved to appendix or omitted)
%%   - Gossip section REMOVED from main body
%%   - New §5.X: Scaling law IEE ∝ n^α characterised empirically
%%   - Abstract reframed for statistical-physics audience
%%   - Engineering framing de-emphasised; physics narrative throughout
%%
%% Compile: pdflatex ELAPSE_PhysicaA.tex
%%          bibtex   ELAPSE_PhysicaA
%%          pdflatex ELAPSE_PhysicaA.tex  (x2)
%%
\documentclass[review,12pt]{elsarticle}

\usepackage{amsmath,amssymb,amsfonts}
\usepackage{amsthm}
\usepackage{graphicx}
\usepackage{booktabs}
\usepackage{hyperref}
\usepackage{xcolor}
\usepackage{url}
\usepackage{natbib}
\usepackage{multirow}
\usepackage{caption}
\usepackage{subcaption}
\usepackage{pifont}

%% For new simulation figures: use \includefigure to include with fallback.
%% If the figure file does not exist, a placeholder text is shown.
\newcommand{\includefigure}[2][width=\textwidth]{%
  \includegraphics[#1]{#2}%
}

%% Custom commands
\newcommand{\R}{\mathbb{R}}
\newcommand{\E}{\mathbb{E}}
\newcommand{\ind}{\mathbf{1}}       % indicator function
\newcommand{\bfx}{\mathbf{x}}
\newcommand{\bfw}{\mathbf{w}}
\newcommand{\bfv}{\mathbf{v}}
\newcommand{\IEE}{\mathrm{IEE}}
\newcommand{\Hc}{H_c}
\newcommand{\Hmax}{H_{\max}}
\newcommand{\lambdat}{\lambda_2}

\newtheorem{theorem}{Theorem}
\newtheorem{lemma}{Lemma}
\newtheorem{observation}{Observation}
\newtheorem{proposition}{Proposition}
\newtheorem{conjecture}{Empirical Conjecture}
\newtheorem{corollary}{Corollary}
\newtheorem{remark}{Remark}
\newtheorem{definition}{Definition}

\graphicspath{{../../output/figures/}}

\journal{Physica A: Statistical Mechanics and its Applications}

\begin{document}

\begin{frontmatter}

\title{ELAPSE: Entropy-Linked Autonomous Propagation and Self-Erasure ---
       A Multi-Domain Ensemble Framework for Active Data Containment
       in Decentralised Networks}

\author[rv]{Y.~Sailaja}
\ead{sailaja@rvce.edu.in}

\author[rv]{Chhavi Pareek}
\author[rv]{Hitarth Mehra}
\author[rv]{Harsh Patel}

\address[rv]{Department of Mathematics / Department of Computer Science \&
             Engineering, RV College of Engineering, Bangalore 560059, India}

\begin{abstract}
We study information spreading and active containment on complex networks
through the lens of statistical mechanics.  A data payload diffusing on a
network drives its spatial Shannon entropy $H(t) = -\sum_i p_i \ln p_i$
from a localised initial state toward maximum entropy $\Hmax = \ln n$,
analogous to a thermal relaxation.  We show that this trajectory exhibits a sharp \emph{crossover regime}
near a critical threshold $\Hc^*$:
below $\Hc^*$, entropy-triggered deletion fires before the data infects the
bulk (subcritical); above $\Hc^*$, the system saturates before deletion acts
(supercritical).  The crossover width $\sim 1/(2\alpha\lambda_2 T)$ is
controlled by the network's algebraic connectivity $\lambda_2$.
Finite-size scaling analysis shows that $\Hc^*$ shifts with network size,
indicating a crossover rather than a sharp universal phase transition.

We introduce ELAPSE (Entropy-Linked Autonomous Propagation and Self-Erasure),
a comparative framework for five physics-motivated deletion mechanisms ---
diffusion (M1), epidemic SIR (M2), stochastic first-passage (M3), cooperative
binding (M4), and local cascade pressure (M5), fused via an
entropy-regularised ensemble (M6), with: (i) well-posedness via
Yamada--Watanabe; (ii) a spectral IEE bound
$\IEE \leq \ln(n)\,M_0\,[\tau(\lambda_2) + O(\mu^{-1})]$;
(iii) IEE scaling IEE~$\sim n^\alpha$ with $\alpha \approx 1.2$--$1.4$
across Erd\H{o}s-R\'enyi, Barab\'asi-Albert, and Watts-Strogatz topologies;
(iv) exact convexity of the linearised objective.
Statistical evaluation over $N_{\mathrm{test}}=20$ seeds at $n \in \{50,
100, 200, 500, 1000\}$ establishes $\geq 9\times$ IEE reduction over the
uncontrolled baseline.  The principal quantitative finding is that the
local cascade pressure mechanism (M5) achieves the lowest Time-Integrated
Entropy Exposure across all three canonical network topologies, consistent
with the known role of local neighbourhood dynamics in complex contagion.
The ensemble (M6) is competitive with M5 under topology-matched training
and serves as an adaptive mechanism-selection procedure when the optimal
mechanism is not known \emph{a priori}.
\end{abstract}

\begin{keyword}
Complex networks \sep Information diffusion \sep Shannon entropy \sep
Phase transition \sep Percolation \sep Stochastic differential equations
\end{keyword}

\end{frontmatter}

%% ─── 1. INTRODUCTION ───────────────────────────────────────────────────────
\section{Introduction}
\label{sec:intro}

\subsection{Information Diffusion on Complex Networks}

The spreading of information, pathogens, and influence across complex networks
is one of the central problems of statistical mechanics on graphs.
Foundational work by Barab\'asi and Albert~\cite{barabasi1999emergence} on
scale-free networks and by Watts and Strogatz~\cite{watts1998collective} on
small-world topologies established that heterogeneous degree distributions and
clustering structure fundamentally alter diffusion thresholds.  Pastor-Satorras
and Vespignani~\cite{pastor2001epidemic} demonstrated that epidemic spreading on
scale-free networks is threshold-free in the infinite-size limit, a result with
deep implications for any diffusion-containment strategy.  The comprehensive
reviews by Albert and Barab\'asi~\cite{albert2002statistical} and Boccaletti
et al.~\cite{boccaletti2006complex} established complex network science as a
mature field within statistical physics, characterising how the interplay of
topology and dynamics produces emergent collective behaviour.

The \emph{controlled} version of diffusion (halting or reversing a spreading
process before it saturates the network) has received comparatively less
attention.  Classical epidemic-control literature focuses on vaccination
thresholds and quarantine timing~\cite{newman2003structure,
dorogovtsev2008critical}, but these strategies target the spreading agent rather
than the payload itself.  In information-diffusion settings, where the
``payload'' is a data object rather than a pathogen, the relevant control
objective is to minimise the total network-wide exposure to the data before
it is irretrievably dispersed.  This problem arises concretely whenever a data
payload in a decentralised peer-to-peer (P2P) network must be autonomously
contained after initial release: once a payload has diffused across hundreds of
pseudonymous P2P nodes, passive containment strategies are ineffective.

\subsection{Entropy as an Order Parameter for Diffusion Control}

The Shannon entropy $H(t) = -\sum_i p_i \ln p_i$ of the spatial concentration
distribution $\{p_i(t)\}$ provides a natural order parameter for the degree of
information spread~\cite{zhao2024multi}.  When data is concentrated on few nodes,
$H(t) \approx 0$ (ordered phase); when data has spread uniformly across the
network, $H(t) \to \Hmax = \ln n$ (disordered phase).  The temporal trajectory
$H(t)$ encodes the collective spreading dynamics in a single scalar observable,
independent of the specific network realisation, analogous to a
thermodynamic observable that integrates over microscopic degrees of freedom.

This entropy trajectory is not merely a passive descriptor but can serve as a
\emph{feedback signal} for active control.  If a self-erasure mechanism is
triggered when $H(t)$ crosses a critical threshold $\Hc$, then the deletion
event is conditioned on the actual spatial state of the diffusion process
rather than an elapsed wall-clock time.  The resulting \emph{entropy clock}
is topology-aware: on a dense Erd\H{o}s-R\'enyi graph with high algebraic
connectivity, entropy rises rapidly and the trigger fires early; on a sparse
Watts-Strogatz ring lattice with low algebraic connectivity, entropy rises
slowly and deletion waits until meaningful spreading has occurred.

\subsection{Limitations of Chronological Containment}

The existing active-containment literature (most influentially the
Vanish~\cite{geambasu2009vanish} and FADE~\cite{tang2010fade} systems) relies
on a shared architectural assumption: a \emph{chronological decay timer} governs
data lifetime. Vanish overwrites Distributed Hash Table (DHT) shards on a fixed
schedule (typically 8~hours), so that the decryption key is irretrievably lost
after a wall-clock timeout. FADE extends this with file-assured deletion via
policy-based key management.

The fundamental limitation of this paradigm is its independence from the actual
network topology. On a scale-free (Barab\'asi-Albert) graph, the hub-driven
spreading mechanism can saturate a majority of nodes within a small fraction of
any fixed timeout~\cite{barabasi1999emergence,pastor2015epidemic}. A data payload
that has already reached $10^3$ nodes before the 8-hour timer fires may be
practically irrecoverable regardless of key deletion. Conversely, on a
sparse small-world (Watts-Strogatz) graph, the same timer may fire when the
data has barely left the source neighbourhood, over-triggering deletion.

\subsection{The Entropy-Clock Approach}

We propose that data lifetime should be a \emph{dynamic function of the
network state}. Specifically, ELAPSE uses the Shannon entropy of the data's
spatial distribution as its deletion clock:
\begin{equation}
  H(t) = -\sum_{i=1}^{n} p_i(t) \ln p_i(t), \quad
  p_i(t) = \frac{x_i(t)}{\sum_j x_j(t)},
  \label{eq:entropy}
\end{equation}
where $x_i(t) \geq 0$ is the data concentration at node $i$ at time $t$.
When $H(t)$ is low, data is concentrated on few nodes (low privacy exposure);
when $H(t)$ approaches $\Hmax = \ln n$, data has spread uniformly (maximal
exposure). Self-erasure is triggered when $H(t)$ crosses a critical
threshold $\Hc$, replacing the chronological clock with a
\emph{spatial-topological clock}.

\subsection{Contributions}

This paper makes the following contributions in complex network statistical
mechanics and controlled diffusion on graphs:
\begin{enumerate}
  \item \textbf{Entropy-driven diffusion framework and mechanism ranking}:
        We introduce ELAPSE (Entropy-Linked Autonomous Propagation and
        Self-Erasure), a comparative framework for five physics-motivated
        deletion mechanisms drawn from diffusion PDEs, epidemic compartmental
        models, stochastic first-passage theory, cooperative binding functions,
        and social cascade dynamics, evaluated under a unified IEE objective
        (Sections~\ref{sec:mechanisms}--\ref{sec:ensemble}).  The statistical
        evaluation establishes social cascade pressure (M5) as the
        dominant mechanism: it achieves the lowest IEE on all three
        network topologies, a finding that is accessible only because
        ELAPSE provides a common entropy-clock framework within which
        mechanisms from disparate physical disciplines can be compared
        on equal footing.
  \item \textbf{Rigorous statistical mechanics results}
        (Section~\ref{sec:theory}): Five theoretical results establishing
        ELAPSE as a well-defined dynamical system with computable bounds:
        (i)~pathwise-unique non-negative strong solution on $[0,T]$
        (Theorem~\ref{thm:wellposed}, via Yamada--Watanabe);
        (ii)~a closed-form spectral IEE upper bound in terms of network size $n$,
        initial mass $M_0$, and deletion intensity
        (Theorem~\ref{thm:iee_bound});
        (iii)~a topology-refined bound incorporating the algebraic connectivity
        $\lambda_2$ (Fiedler value), giving the spectral trigger time
        $\tau(\lambda_2)$ explicitly (Corollary~\ref{cor:lambda2});
        (iv)~a sharp crossover regime near the critical entropy
        threshold $\Hc^*$ separating subcritical (effective containment)
        and supercritical (bulk infection) regimes, with crossover width
        controlled by $\lambda_2$ (Section~\ref{sec:phase_transition});
        (v)~exact convexity of the linearised IEE objective
        (Theorem~\ref{thm:lin_convex}), with the full nonlinear surface
        exhibiting a 14.1\% Jensen-violation rate.
  \item \textbf{Topology-dependent scaling laws}: IEE $\sim n^\alpha$ with
        exponents $\alpha \in [1.12, 1.40]$ that depend systematically on
        network topology and algebraic connectivity, consistent with the
        spectral trigger time prediction $\tau \sim \lambda_2^{-1}$
        (Section~\ref{sec:scaling_laws}).
  \item \textbf{Statistical evaluation}: All IEE values reported as
        mean $\pm$ 95\% CI over $N_{\mathrm{test}} = 20$ out-of-sample seeds,
        three topologies (ER/BA/WS), and four network sizes
        $n \in \{50, 100, 200, 500\}$ (Section~\ref{sec:results}).
  \item \textbf{Real-world and distributed extensions}: ELAPSE is validated
        on Stanford SNAP Gnutella P2P subgraphs (Section~\ref{sec:snap}),
        and a gossip-protocol distributed entropy estimator is proposed and
        characterised in terms of minimum hop count $k^* = O(1/\lambda_2)$
        (Appendix~\ref{sec:gossip}).
\end{enumerate}

\subsection{Paper Organisation}

Section~\ref{sec:related} reviews the literature and positions ELAPSE.
Section~\ref{sec:mechanisms} defines the five core mechanisms.
Section~\ref{sec:ensemble} presents the ensemble construction and five
theoretical results (Theorems~\ref{thm:wellposed}--\ref{thm:lin_convex},
Corollaries~\ref{cor:lambda2}--\ref{cor:unique},
Propositions~\ref{prop:privacy}--\ref{prop:vs_timer}).
Section~\ref{sec:results} provides the main statistical evaluation.
Section~\ref{sec:snap} reports real-world validation.
Section~\ref{sec:adversarial} analyses adversarial evasion.
Section~\ref{sec:conclusions} concludes.
Appendix~\ref{sec:gossip} covers distributed entropy estimation
(gossip protocol); Appendix~\ref{sec:m3} provides a mechanistic
analysis and fix for M3's underperformance.

%% ─── 2. RELATED WORK ───────────────────────────────────────────────────────
\section{Related Work}
\label{sec:related}

\subsection{Information Diffusion and Entropy on Graphs}

The statistical mechanics of information spreading on complex networks is
a central topic in Physica A and related journals.
The comprehensive review by Albert and Barab\'asi~\cite{albert2002statistical}
established that heterogeneous degree distributions fundamentally alter
spreading threshold behaviour: scale-free networks lack an epidemic threshold
in the infinite-size limit~\cite{pastor2001epidemic}, implying that diffusion
processes spread to a non-zero fraction of the network for any non-zero
transmission rate.  The interplay of network topology, spreading dynamics,
and emergent phase transitions is reviewed in
Boccaletti et al.~\cite{boccaletti2006complex} and
Newman~\cite{newman2003structure}.

The particular role of Shannon entropy $H(t) = -\sum_i p_i \ln p_i$ as a
macroscopic order parameter for diffusion state has gained prominence through
Zhao et al.~\cite{zhao2024multi}, who demonstrated that entropy-spreading
metrics accurately proxy information volatility across heterogeneous degree
distributions, establishing $H(t)$ as a reliable outbreak-detection signal.
Epidemic spreading on complex networks, from SIR to SIS and beyond,
is surveyed in Pastor-Satorras et al.~\cite{pastor2015epidemic}, who
establish that node degree correlates with epidemic exposure, a result
connected to our per-topology IEE scaling laws (Section~\ref{sec:scaling_laws}).
Moreno et al.~\cite{moreno2002epidemic} further characterise how heterogeneous
degree distributions accelerate epidemic outbreaks in scale-free topologies,
consistent with ELAPSE's observation that BA networks (higher-degree hubs)
demand faster entropy-triggered deletion.

Parallel advances in continuous graph diffusion, originating in spectral graph
signal processing~\cite{shuman2013emerging} and extended through neural
PDE-on-graph formulations~\cite{chamberlain2021grand,eliasof2024graph,fan2024graph},
confirm that partial differential equations on graph topologies accurately
capture non-linear propagation boundaries.  Universal resilience
patterns~\cite{gao2016universal} characterise scale-stress behaviour under
perturbation, providing a complementary macroscopic lens to the entropy
trajectories studied here.

\subsection{Percolation, Phase Transitions, and Self-Organised Criticality}

Percolation theory on complex graphs provides the natural framework for
understanding threshold phenomena in network spreading.
Callaway et al.~\cite{callaway2000network} established that random-graph
percolation transitions can be mapped onto epidemic threshold conditions,
and Dorogovtsev et al.~\cite{dorogovtsev2008critical} provide a unified
review of critical phenomena on complex networks.  The critical threshold
$\Hc^*$ in ELAPSE (Section~\ref{sec:phase_transition}) plays a role
analogous to the bond-percolation threshold $p_c$: below $\Hc^*$, deletion
fires before data infects the bulk (subcritical); above $\Hc^*$, the
network saturates before deletion can act (supercritical).
Newman et al.~\cite{newman2001random} show that degree-heterogeneity
shifts this transition point, consistent with our finding that $\Hc^*$
varies across ER, BA, and WS topologies through the algebraic connectivity
$\lambda_2$.  Self-organised criticality~\cite{vespignani2012modelling}
provides a complementary perspective: ELAPSE self-tunes the deletion onset
to the network's current entropy state rather than imposing an external
threshold.

\subsection{Static and Active Data Containment}

Early containment schemes (Vanish~\cite{geambasu2009vanish},
FADE~\cite{tang2010fade}, and the Ephemerizer~\cite{perlman2005ephemerizer})
rely on \emph{chronological} timers that are independent of network
topology.  Wolchok et al.~\cite{wolchok2010defeating} demonstrated Sybil-attack
vulnerability in DHT-based schemes.  Rondeau and Temple~\cite{rondeau2021secure}
survey secure deletion approaches and confirm that time-blind mechanisms are
structurally insufficient for dynamic spreading environments.
The fundamental gap is that a fixed timeout applied to a scale-free
network~\cite{albert2002statistical} either fires too late
(hub-driven spreading saturates the network before the timer) or too
early (sparse periphery barely reached when deletion triggers).

\subsection{Algorithmic Forgetting and Entropy-Based Privacy}

Complementary to network-level containment, the machine-learning literature
has developed \emph{algorithmic forgetting} (machine
unlearning~\cite{bourtoule2021machine,xu2023machine}): removing a training
sample's influence from a model.  Both problems share the use of
entropy-based metrics as core sensitivity measures.  The connection to
differential privacy~\cite{dwork2014algorithmic} (which bounds inference
leakage from a fixed dataset) is formal but distinct: ELAPSE bounds
the \emph{spatial physical exposure} of a diffusing payload, not inference
leakage.  Both frameworks encode privacy via an exponential decay parameter;
Proposition~\ref{prop:privacy} shows that ELAPSE achieves a mass-retention
guarantee of the same functional form as $\varepsilon$-differential privacy.

\subsection{Control of Diffusion Processes on Graphs}

ELAPSE is related to, but distinct from, the emerging literature on
\emph{controlled diffusion on networks}.
Mean-field game (MFG) frameworks~\cite{carmona2018probabilistic,achdou2020mean}
study populations of McKean-–Vlasov-type agents interacting through their
empirical distribution, deriving feedback control laws from coupled
Hamilton-–Jacobi-–Bellman and Fokker-–Planck equations.
In MFG language, ELAPSE is a \emph{social planner} problem: we seek a
mortality schedule $\mu\Delta(t;\bfw)$ that minimises a population-level cost
(IEE) subject to the coupled McKean-–Vlasov SDE \eqref{eq:egdm}--\eqref{eq:ensemble}.
However, ELAPSE differs from classical MFG in three ways: (i)~topology enters
via the Laplacian $L$ rather than a mean-field coupling term; (ii)~the control
$\Delta(t;\bfw)$ is an entropy-dependent functional of the empirical
distribution; and (iii)~we optimise over a finite ensemble of pre-specified
mechanisms rather than solving a PDE.
Online feedback control on network flow
manifolds~\cite{bianchin2021online} addresses a related setting but focuses on
power systems rather than information-theoretic objectives.
ELAPSE contributes to this space with the first entropy-functional IEE bound
(Theorem~\ref{thm:iee_bound}) and spectral trigger time
(Corollary~\ref{cor:lambda2}) for controlled diffusion on heterogeneous P2P graphs.

\subsection{Positioning of ELAPSE}

Table~\ref{tab:positioning} summarises how ELAPSE differs from the three
dominant paradigms for data-lifetime management.

\begin{table}[htbp]
  \caption{Comparison of data-deletion paradigms.  ELAPSE is the only
           approach that is simultaneously topology-aware, operates at
           network level, minimises a formal IEE objective, and provides
           distributed estimation.}
  \label{tab:positioning}
  \centering
  \small
  \renewcommand{\arraystretch}{1.15}
  \begin{tabular}{lcccc}
    \toprule
    Paradigm & Topology- & Network- & Formal & Distributed \\
             & aware     & level    & IEE obj. & estimation \\
    \midrule
    Cryptographic (Vanish, FADE)~\cite{geambasu2009vanish,tang2010fade}
              & \ding{55} & \ding{55} & \ding{55} & \ding{55} \\
    Chron.\ timer (M7, \S\ref{sec:m7_comparison})
              & \ding{55} & \ding{55} & \ding{55} & \ding{55} \\
    Machine unlearning~\cite{bourtoule2021machine,xu2023machine}
              & \ding{55} & \ding{55} & \ding{55} & \ding{55} \\
    Differential privacy~\cite{dwork2014algorithmic}
              & \ding{55} & \ding{55} & \ding{55} & Partial \\
    \textbf{ELAPSE (ours)}
              & \ding{51} & \ding{51} & \ding{51} & \ding{51} \\
    \bottomrule
  \end{tabular}
\end{table}

ELAPSE occupies a unique position in the landscape of controlled diffusion
on complex networks: it is the first framework that
(i)~replaces the chronological erasure clock with a spatial entropy trigger,
establishing a direct link between network topology (via $\lambda_2$) and
the deletion onset time;
(ii)~constructs a regularised ensemble across five distinct physical mechanisms,
each motivated by an independent branch of statistical physics or mathematical
biology;
(iii)~provides rigorous theoretical guarantees
(Theorems~\ref{thm:wellposed}--\ref{thm:lin_convex}) including a spectral IEE
bound that explicitly depends on the Fiedler value of the network Laplacian;
(iv)~identifies a sharp crossover regime in the entropy threshold parameter
space, with crossover width $\Delta H_c \sim 1/(2\alpha\lambda_2 T)$ controlled
by the Fiedler value (Section~\ref{sec:phase_transition});
and (v)~characterises the minimum gossip hop count $k^*$ as a function of
algebraic connectivity (Proposition~\ref{prop:gossip_lambda}).
Prior ensemble approaches for network spreading
(e.g., competing epidemics~\cite{karrer2011competing}; social
cascades~\cite{guilbeault2018complex}) address spreading \emph{dynamics}
but not active containment.  MFG approaches~\cite{carmona2018probabilistic}
study population-level control but via PDE solutions rather than mechanism
ensembles.

%% ─── 3. FIVE CORE MECHANISMS ───────────────────────────────────────────────
\section{The Five Core Mechanisms}
\label{sec:mechanisms}

\subsection{Network and Notation}

Let $G = (V, E)$ be an undirected connected graph with $|V| = n$ nodes. The
graph Laplacian $L = D - A$ (where $D$ is the degree matrix and $A$ the
adjacency matrix) and the Fiedler value $\lambdat$ (second-smallest eigenvalue
of $L$, the algebraic connectivity~\cite{chung1997spectral}) play central roles.
The data concentration at node $i$ at time $t$ is $x_i(t) \geq 0$, with units
of \emph{data mass per node}.  Initial conditions are $x_i(0) = u_i$ where
$u_i \sim \mathrm{Uniform}(0.5, 1.0)$ for a randomly chosen 10\% of nodes
(the initial data sources) and $x_i(0) = 0$ otherwise; total initial mass
$M(0) = \sum_i x_i(0) \in (0, n/10]$.  All IEE values therefore have units
of [nats $\times$ data-mass $\times$ time], and scale with both $n$ and $M(0)$.
Spatial privacy exposure is measured by $H(t)$ as defined in Eq.~\eqref{eq:entropy},
with $\Hmax = \ln n$ and a critical threshold $\Hc = 0.65 \Hmax$ used throughout.

Each of the five mechanisms defines a continuous \emph{deletion vote}
$v_k(t) \in [0,1]$, measuring urgency of erasure. The composite mortality
pressure $\Delta(t) = \sum_k w_k v_k(t)$ drives data removal.

\subsection{M1: Entropy-Gated Diffusion Model (EGDM)}

The physical baseline treats data as a field propagating against a
graph-Laplacian operator with an entropy-triggered global mortality term:
\begin{equation}
  \frac{d\bfx}{dt} = -\alpha L\bfx - \mu\,\sigma(H(t))\,\bfx
                     + \mathbf{s}(t) + dW_t,
  \label{eq:egdm}
\end{equation}
where $\alpha = 0.3$ is the diffusion coefficient, $\mu = 1.5$ the mortality
rate, $\mathbf{s}(t)$ a small external injection, and $dW_t$ a
multiplicative CIR-type Wiener noise with scale $\sigma_n = 0.015$.
The polynomial vote is:
\begin{equation}
  v_1(t) = \left[\max\!\left(0,\;\frac{H(t) - \Hc}{\Hmax - \Hc}\right)\right]^{\!\beta},
  \label{eq:v1}
\end{equation}
with sharpness $\beta = 2$. Near-zero below $\Hc$, the vote rises as a
smooth power-law as data spreads. Strong Lipschitz continuity of the
deterministic flow ensures existence and uniqueness of strong solutions
up to the erasure threshold~\cite{oksendal2003stochastic}.

\subsection{M2: Epidemic-Entropy (Compartmental SIR)}

Epidemiology models data spread as an infectious load. Nodes transition
through Susceptible-Infected-Recovered compartments, capturing local
neighbour-to-neighbour spreading independently of global diffusion:
\begin{equation}
  \frac{dI_i}{dt} = \beta_{\mathrm{sir}}\,S_i \sum_j A_{ij} I_j
                   - \gamma I_i - \mu\,\sigma(H(t))\,I_i,
  \label{eq:sir}
\end{equation}
where $\beta_{\mathrm{sir}} = 0.4$, $\gamma = 0.15$. The vote is the
current infected fraction: $v_2(t) = n^{-1}\!\sum_i I_i(t)$. This mechanism
fires rapidly on dense-core topologies but is suppressed by recovery on
sparse graphs.

\subsection{M3: Stochastic-Finance (Ornstein-Uhlenbeck First-Passage)}

Drawing from barrier-option pricing~\cite{black1973pricing,uhlenbeck1930theory},
information exposure risk is modelled as a mean-reverting stochastic process
approaching a ruin boundary. The state evolves as:
\begin{equation}
  d\bfx = \theta(\mu_{ou} - \bfx)\,dt + \sigma_{ou}\,dW_t
         - \mu\,P(\tau \leq t)\,\bfx\,dt,
  \label{eq:finance}
\end{equation}
with $\theta = 0.3$, $\mu_{ou} = 0.1$, $\sigma_{ou} = 0.04$. The vote
is the first-passage probability:
\begin{equation}
  v_3(t) = 1 - \exp\!\left(-\lambda_{\mathrm{fp}}\,t\right), \quad
  \lambda_{\mathrm{fp}} = \frac{\alpha\lambdat}{\Hmax - \Hc},
  \label{eq:v3}
\end{equation}
encoding the spectral geometry of the network through the Fiedler
value $\lambdat$. As analysed in Section~\ref{sec:m3}, the static use
of $\lambdat$ limits M3's responsiveness to fast BA hub dynamics.

\subsection{M4: Biological Cooperative Binding (Hill Function)}

In systems biology, cooperative binding transitions are modelled by the
Hill function~\cite{alon2019systems}:
\begin{equation}
  v_4(t) = \frac{H(t)^{n_h}}{H_c^{n_h} + H(t)^{n_h}},
  \label{eq:hill}
\end{equation}
with Hill coefficient $n_h = 4$. This produces a smoother, more
Lyapunov-stable asymptotic transition than the polynomial $v_1$: it is
near zero for $H \ll \Hc$, exactly $\tfrac{1}{2}$ at $H = \Hc$, and
near one for $H \gg \Hc$. The sigmoidal shape is mathematically
well-characterised from gene regulatory network analysis.

\subsection{M5: Sociological Cascade Pressure}

Sociological cascade models parameterise agents that delete based on
localised sub-graph density~\cite{watts2002simple,guilbeault2018complex}.

\begin{definition}[Indicator notation]
\label{def:indicator}
For any predicate $P$, let $\ind[P]$ denote the indicator function:
$\ind[P] = 1$ if $P$ is true, $\ind[P] = 0$ otherwise.
Specifically, for a threshold $\kappa$ and deletion threshold $\epsilon$:
$\ind[x_j < \epsilon]$ equals 1 if node $j$ has effectively deleted its
payload (concentration below $\epsilon = 0.05$), and 0 otherwise;
$\ind[\cdot > \kappa]$ equals 1 if the argument exceeds the cascade
threshold $\kappa = 0.3$.
\end{definition}

The cascade vote combines global and local deletion pressure:
\begin{equation}
  v_5(t) = \tfrac{1}{2}\sigma\!\bigl(H(t)\bigr)
           + \tfrac{1}{2n}\sum_{i=1}^{n}
  \ind\!\!\left[
    \frac{\sum_{j \in N(i)}\ind[x_j < \epsilon]}{\deg(i)}
    > \kappa
  \right],
  \label{eq:social}
\end{equation}
where $N(i)$ is the neighbourhood of node $i$ and $\deg(i)$ its degree.
A node responds to deletion when the global entropy exceeds $\Hc$
\emph{or} when more than $\kappa$ of its neighbours have already deleted.

%% ─── 4. ENSEMBLE CONSTRUCTION ─────────────────────────────────────────────
\section{The ELAPSE Ensemble}
\label{sec:ensemble}

\subsection{Ensemble Architecture}

The M6 ELAPSE engine fuses the five votes into a terminal deletion pressure:
\begin{equation}
  \Delta(t) = \bfw^\top \bfv(t), \quad
  \bfw \in \mathcal{W} = \!\left\{\bfw \in \R^5_{\geq 0} : \|\bfw\|_1 = 1\right\},
  \label{eq:ensemble}
\end{equation}
where $\bfv(t) = [v_1(t),\ldots,v_5(t)]^\top$.
The Time-Integrated Entropy Exposure (IEE) is:
\begin{equation}
  \IEE(\bfw) = \int_0^T H\!\bigl(\mathbf{p}(t;\bfw)\bigr)\cdot M(t;\bfw)\,dt,
  \label{eq:iee}
\end{equation}
where $M(t;\bfw) = \sum_i x_i(t)$ is the surviving data mass. A lower
IEE indicates better privacy protection.

% iee_axiomatic_justification.tex
% Axiomatic justification of IEE as a privacy metric.
% Insert after the IEE definition in Section 4 (Ensemble Architecture).

\begin{remark}[Axiomatic justification of IEE]
\label{rem:iee_axiomatic}
The Time-Integrated Entropy Exposure objective $\IEE(\bfw) = \int_0^T
H(\mathbf{p}(t;\bfw))\cdot M(t;\bfw)\,dt$ is the unique functional
satisfying three natural privacy axioms:

\begin{enumerate}[label=(\roman*)]
  \item \textbf{Mass sensitivity.}  A data deletion mechanism that removes
    mass quickly should be rewarded: if $M(t;\bfw_1) \le M(t;\bfw_2)$
    pointwise on $[0,T]$ and $H(\mathbf{p}(t;\bfw_1)) = H(\mathbf{p}(t;\bfw_2))$,
    then $\IEE(\bfw_1) \le \IEE(\bfw_2)$.
    The product $H \cdot M$ captures this: equal entropy spread over less
    surviving mass means less total exposure.

  \item \textbf{Entropy sensitivity.}  A mechanism that achieves rapid
    localisation of the data distribution (low entropy) while retaining
    mass should also be rewarded.  If $H(\mathbf{p}(t;\bfw_1)) \le
    H(\mathbf{p}(t;\bfw_2))$ pointwise and $M(t;\bfw_1) = M(t;\bfw_2)$,
    then $\IEE(\bfw_1) \le \IEE(\bfw_2)$.
    This prevents the degenerate solution of retaining all data in one node
    (low entropy but high mass exposure) from being penalised inconsistently.

  \item \textbf{Temporal accumulation.}  Privacy exposure is additive in
    time: a protocol that leaks data over a longer window is strictly worse
    than one that resolves the same snapshot exposure faster.  The integral
    $\int_0^T$ implements this.
\end{enumerate}

Any separable functional $\int_0^T f(H, M)\,dt$ satisfying (i)--(iii) with
$f$ increasing in both arguments and $f(0,M) = f(H,0) = 0$ reduces to
$f(H,M) = c\,H\cdot M$ for some constant $c > 0$ under a standard
homogeneity argument analogous to Shannon's axiomatic derivation of
entropy \cite{jaynes1957information}.  Setting $c=1$ gives the IEE.  Thus IEE is not
an ad-hoc combination but the unique canonical measure of time-integrated
privacy exposure under axioms (i)--(iii).
\end{remark}


\subsection{Linear Vote Fusion: Derivation and Alternatives}
\label{sec:fusion_justification}

The linear combination $\Delta(t) = \bfw^\top\bfv(t)$ is the natural
parametrisation of additive-mortality control in the ELAPSE SDE.

\begin{proposition}[Linear fusion is the canonical implementation of additive-mortality control]
\label{prop:fusion}
In the ELAPSE SDE, mortality enters as $-\mu\Delta(t;\bfw)\,x_i$ with
$\Delta(t;\bfw) = \bfw^\top\bfv(t)$, so that mechanism $k$ contributes
deletion rate $\mu w_k v_k(t)$ independently of other mechanisms.
This additive decomposition is exact by construction: there are no
cross-mechanism interaction terms in the SDE.  Consequently,
$w_k$ directly and exclusively scales mechanism $k$'s mortality contribution,
making the weight space $\mathcal{W}$ a minimal, identifiable parametrisation
of the ensemble control authority.  Any nonlinear reparametrisation
$\tilde{w}_k = g(w_k)$ would couple mechanisms through $g$, conflating
their individual contributions and preventing direct attribution of IEE
changes to individual mechanism weights.
\end{proposition}

\begin{remark}[Scope of the proposition]
Proposition~\ref{prop:fusion} justifies linear fusion as \emph{structurally
correct} for additive SDE mortality, not as universally optimal among all
conceivable fusion rules.  Whether nonlinear fusion (e.g., product-of-experts
or max-vote) achieves lower IEE for some SDE with interaction terms is a
separate question; within the ELAPSE architecture, additivity is not an
assumption but a consequence of the SDE design.
\end{remark}

\paragraph{Comparison with nonlinear fusion rules.}
Two natural alternatives are the \emph{product-of-experts}
$\Delta_{\rm PoE}(t) = \prod_k v_k(t)^{w_k}$
(geometric mean) and the \emph{maximum-vote}
$\Delta_{\rm max}(t) = \max_k (w_k v_k(t))$.
Product-of-experts requires all mechanisms to simultaneously vote for deletion
(logical AND), which is over-conservative and delays $\Delta$ even when four
of five mechanisms signal high entropy.  Maximum-vote uses only the single
most-confident mechanism, ignoring the diversity of evidence.  The linear
combination provides a smooth interpolation: it fires when the
\emph{weighted majority} of mechanisms vote positively.
Quantitative IEE comparison of these fusion classes is identified as future work,
as it requires independent ensemble optimisation for each class.

\subsection{Entropy-Regularised Optimisation and Maximum Entropy Principle}
\label{sec:max_entropy}

Unconstrained Nelder-Mead on $\mathcal{W}$ collapses to a one-hot vertex
$\mathbf{e}_k$ (choosing the single best mechanism), sacrificing
generalisation across topology types. We add an entropy regulariser:
\begin{equation}
  \min_{\bfw \in \mathcal{W}}\; J(\bfw)
    = \IEE(\bfw) - \lambda_{\mathrm{reg}}\,\mathcal{H}(\bfw),
  \label{eq:objective}
\end{equation}
where $\mathcal{H}(\bfw) = -\sum_k w_k \ln w_k$ is the weight-layer
entropy. Subtracting this term penalises low-entropy (collapsed) weight
distributions while allowing IEE minimisation to proceed. We set
$\lambda_{\mathrm{reg}} = 15$ throughout, selected at the Pareto elbow
between IEE and weight diversity (see Table~\ref{tab:reg}).

\paragraph{Connection to maximum entropy principle.}
The regularised objective \eqref{eq:objective} has a natural
statistical-mechanical interpretation via Jaynes' maximum entropy
principle~\cite{jaynes1957information}: given a constraint on the expected
IEE performance, the maximum-entropy distribution over mechanisms is
precisely $\bfw^* = \arg\min J(\bfw)$.  Equivalently, $\bfw^*$ is the
Gibbs distribution $w_k \propto \exp(-\IEE_k / \lambda_{\mathrm{reg}})$
in the limit where each mechanism's marginal IEE $\IEE_k$ is decoupled.
The entropy regulariser thus acts as a \emph{temperature} parameter:
high $\lambda_{\mathrm{reg}}$ enforces a uniform (high-temperature) weight
distribution that maximises mechanism diversity at the cost of IEE
performance; low $\lambda_{\mathrm{reg}} \to 0$ collapses to the zero-temperature
ground state (single best mechanism).  The Pareto elbow at $\lambda_{\mathrm{reg}} = 15$
(Table~\ref{tab:reg}) identifies the critical temperature below which
ensemble diversity begins to collapse.

\subsection{Theoretical Analysis}
\label{sec:theory}

We establish three main results: well-posedness of the SDE system, a
closed-form spectral upper bound on IEE, and rigorous convexity of the
linearised IEE objective.

\noindent\emph{Physical intuition}: The ELAPSE system is a
CIR-type stochastic process (Cox-Ingersoll-Ross), well-studied in
mathematical finance and population dynamics.  Non-negativity follows from
the vanishing diffusion coefficient at the zero boundary (a ``reflecting
barrier''), and uniqueness follows from the H\"{o}lder-$\frac{1}{2}$
condition on the noise amplitude.

\begin{theorem}[Well-posedness of ELAPSE]
\label{thm:wellposed}
Let $G$ be connected and $\bfx(0) \geq 0$ componentwise.
The ELAPSE SDE system \eqref{eq:egdm}--\eqref{eq:ensemble} has a
pathwise-unique, non-negative strong solution $\bfx(t) \geq \mathbf{0}$
a.s.\ on $[0,T]$ for any finite $T > 0$.
\end{theorem}

\begin{proof}
We establish the four parts in \emph{strict logical order}, so that each step
uses only results already proved (eliminating the circular dependency present
in earlier drafts where global existence was invoked before non-negativity).

\medskip
\noindent\textbf{Step A: Local existence.}
On any compact $K \subset \mathbb{R}^n_+$, the drift
$b_i(\bfx,t) = -\alpha(L\bfx)_i - \mu\Delta(t;\bfw)x_i + s_i(t)$
is Lipschitz in $\bfx$ (the Laplacian term is linear; $-\mu\Delta x_i$ is
Lipschitz with constant $\mu$; $s_i \geq 0$ is constant), and the diffusion
$\sigma_i(\bfx) = \sigma_n\sqrt{x_i}$ satisfies the Yamada--Watanabe
$\tfrac{1}{2}$-H\"older condition $|\sigma_i(x)-\sigma_i(y)| \leq
\sigma_n\sqrt{|x_i-y_i|}$.  By the SDE Picard--Lindel\"of theorem
\cite[Thm.~5.2.1]{oksendal2003stochastic}, there exists a unique local strong
solution on $[0,\tau_K)$ where $\tau_K = \inf\{t\!:\bfx(t)\notin K\}$.

\medskip
\noindent\textbf{Step B: Non-negativity on $[0,\tau_K)$, \emph{without} global
existence.}
Fix component $i$. At the boundary $x_i = 0$, the diffusion coefficient
$\sigma_n\sqrt{0}=0$ vanishes exactly, and the drift satisfies
\[
  b_i(\bfx,t)\big|_{x_i=0}
  \;=\; \alpha\!\sum_{j\sim i} A_{ij}\,x_j + s_i \;\geq\; 0,
\]
since $A_{ij}\geq 0$, $x_j \geq 0$ on $[0,\tau_K)$ (componentwise, by the
same argument inductively for each $j$), and $s_i \geq 0$.
By Feller's boundary classification~\cite{feller1951two}, the boundary $x_i=0$
is \emph{entrance} (non-negative drift, vanishing diffusion), meaning the
process cannot reach it from $x_i(0) > 0$ in finite time.
Formally: apply the comparison theorem~\cite[Thm.~IX.3.7]{revuzyor1999} with
the dominating zero-drift process $Y_i$ satisfying
$dY_i = \sigma_n\sqrt{Y_i}\,dW_i$, $Y_i(0) = 0$.
Since $b_i|_{x_i=0} \geq 0 =$ (drift of $Y_i$), the comparison theorem gives
$x_i(t) \geq Y_i(t) = 0$ a.s.\ on $[0,\tau_K)$.
Hence $x_i(t) \geq 0$ for all $i$, established independently of global existence.

\medskip
\noindent\textbf{Step C: Global existence.}
\emph{Now} that non-negativity is established on $[0,\tau_K)$, we have
$M(t) = \sum_i x_i(t) \geq 0$ on $[0,\tau_K)$.
Apply Itô's formula to $V(\bfx) = M(\bfx)$:
\[
  dM = \Bigl(-\mu\Delta(t;\bfw)M + \|{\bf s}\|_1\Bigr)dt
       + \sigma_n\sum_i \sqrt{x_i}\,dW_i,
\]
using $\mathbf{1}^\top L = \mathbf{0}$.  Taking expectations (the stochastic
integral is a square-integrable martingale since $\mathbb{E}[\sqrt{x_i}]
< \infty$ by non-negativity and $\mathbb{E}[M] < \infty$ by induction):
\[
  \frac{d}{dt}\,\mathbb{E}[M(t)] \;\leq\; \|{\bf s}\|_1,
  \qquad\Rightarrow\qquad
  \mathbb{E}[M(t)] \;\leq\; M_0 + \|{\bf s}\|_1 T \;<\; \infty.
\]
By Markov's inequality, $P(M(t) > R) \leq (M_0 + \|{\bf s}\|_1 T)/R \to 0$
as $R \to \infty$, so $P(\tau_K < T) \to 0$ as $K \nearrow \mathbb{R}^n_+$.
The solution exists globally on $[0,T]$.

\medskip
\noindent\textbf{Step D: Pathwise uniqueness via Yamada--Watanabe.}
The diffusion coefficient $\sigma(u) = \sigma_n\sqrt{u}$ satisfies
$|\sigma(u) - \sigma(v)| \leq \sigma_n\sqrt{|u-v|}$ for $u,v \geq 0$.
With $h(r) = \sigma_n\sqrt{r}$: $\int_0^1 h(r)^{-2}\,dr
= \sigma_n^{-2}\int_0^1 r^{-1}\,dr = +\infty$,
satisfying the Yamada--Watanabe H\"older-$\tfrac{1}{2}$ condition
\cite[Thm.~IX.3.5]{revuzyor1999}.
Combined with Steps A--C, the Yamada--Watanabe theorem yields a
pathwise-unique strong solution.
\end{proof}

\begin{remark}[Logical ordering removes circularity]
\label{rem:thm1_ordering}
The order Steps A$\to$B$\to$C$\to$D is essential. Non-negativity (Step~B) is
proved using only the local solution from Step~A, via the Feller boundary
argument that requires neither $M(t)>0$ nor global existence. The
Lyapunov--Gronwall bound (Step~C) then \emph{uses} $M(t)\geq 0$ (proved in B)
to conclude $\mathbb{E}[M(t)] < \infty$. Any proof that invokes $M(t) > 0$
before establishing non-negativity would be circular.
\end{remark}

\begin{lemma}[Positivity of deletion intensity]
\label{lem:delta_pos}
Let $\bfw^* \in \mathrm{int}(\mathcal{W})$ be the ensemble weight vector
guaranteed by Corollary~\ref{cor:unique}.  Define
$\tau^* = \tau(\lambdat)$ from \eqref{eq:tau_lambda}.  Then:
\begin{equation}
  \delta_{\min} \;:=\; \inf_{t \in [\tau^*, T]} \Delta(t;\bfw^*) \;>\; 0.
  \label{eq:delta_min}
\end{equation}
\end{lemma}

\begin{proof}
Corollary~\ref{cor:lambda2} gives $H(t) \geq \Hc$ for all $t \geq \tau^*$.
Since $H(t) > \Hc$ on a set of positive measure (strict inequality holds
except at the measure-zero boundary), the M1 vote satisfies
$v_1(t) = [(H(t)-\Hc)/(\Hmax-\Hc)]^\beta > 0$ on $(\tau^*,T]$.
As $\bfw^* \in \mathrm{int}(\mathcal{W})$ we have $w_1^* > 0$, so
$\Delta(t;\bfw^*) \geq w_1^* v_1(t) > 0$.
Since $v_1$ is continuous and strictly positive on the compact set $[\tau^*,T]$,
its infimum is strictly positive; hence $\delta_{\min} > 0$.
\end{proof}

\noindent\emph{Physical intuition}: Once the entropy clock fires (at time $\tau^*$),
the data mass $M(t)$ decays exponentially, and since $H(t) \leq \ln n$
throughout, the remaining IEE accumulation is bounded by an integral of
an exponential (a one-sided Laplace transform).  The bound separates into
a topology-dependent pre-trigger cost (how long before deletion starts)
and a topology-independent post-trigger decay (how fast deletion proceeds).

\begin{theorem}[IEE Upper Bound for the Stochastic ELAPSE System]
\label{thm:iee_bound}
Let $\delta_{\min} > 0$ be as in Lemma~\ref{lem:delta_pos} and
$\tau^* = \tau(\lambdat)$ from \eqref{eq:tau_lambda}.
Define the Gronwall constants
\begin{equation}
  \gamma := -2\mu\,\delta_{\min} + \sigma_n^2,
  \qquad
  C_1 := (2\|{\bf s}\|_1 + \sigma_n^2)\,M_0,
  \label{eq:gronwall_constants}
\end{equation}
and the second-moment envelope
$\overline{M^2}(t) := (M_0^2 + C_1\,t)\,e^{\gamma t}$.

\medskip
\noindent\textbf{(a) Pre-trigger phase $[0, \tau^*]$:}
$\displaystyle\int_0^{\tau^*} H(t)M(t)\,dt \leq \ln(n)\,M_0\,\tau^*$
(valid for the full stochastic system, since $H \leq \ln n$ and $M \leq M_0$
hold a.s.)

\medskip
\noindent\textbf{(b) Stochastic post-trigger bound:}
For the full ELAPSE SDE \eqref{eq:egdm}--\eqref{eq:ensemble} with $\sigma_n \geq 0$:
\begin{equation}
  \mathbb{E}\bigl[M(t)^2\bigr] \;\leq\; \overline{M^2}(t)
  \qquad\forall\, t \in [\tau^*, T],
  \label{eq:E_M2_bound}
\end{equation}
and consequently
\begin{equation}
  \mathbb{E}[\IEE(\bfw)] \;\leq\;
  \ln(n)\,M_0\,\tau^*
  + \ln(n)\!\int_{\tau^*}^{T} \!\sqrt{\overline{M^2}(t)}\;\mathrm{d}t
  \;=:\; \mathcal{B}(n,M_0,\mu,\delta_{\min},\sigma_n,\|{\bf s}\|_1,\tau^*,T).
  \label{eq:iee_bound_stochastic}
\end{equation}

\medskip
\noindent\textbf{(c) Deterministic limit:}
As $\sigma_n \to 0$, $\gamma \to -2\mu\delta_{\min} < 0$, and
$\mathcal{B}$ reduces to
$\ln(n)\,M_0\!\left[\tau^* + \frac{1-e^{-\mu\delta_{\min}(T-\tau^*)}}{\mu\delta_{\min}}\right]$,
recovering exponential mass decay.
\end{theorem}

\begin{proof}
\emph{Part (a)} is immediate since $H(t) \leq \ln n$ and $M(t) \leq M_0$ a.s.

\emph{Part (b).}
Apply Itô's formula to $M(t)^2$:
\begin{align*}
  d(M^2) &= 2M\,dM + d\langle M\rangle_t \\
  &= 2M\bigl[-\mu\Delta(t)M + \|{\bf s}\|_1\bigr]\,dt
   + \sigma_n^2\!\sum_i x_i\,dt + d\mathcal{M}_t,
\end{align*}
where $d\mathcal{M}_t$ is a local martingale (zero mean) and
$d\langle M\rangle_t = \sigma_n^2\sum_i x_i\,dt = \sigma_n^2 M\,dt$.
Using $\Delta(t) \geq \delta_{\min}$ and $\sum_i x_i = M$:
\[
  \frac{d}{dt}\,\mathbb{E}[M^2]
  \;\leq\; \gamma\,\mathbb{E}[M^2] + 2\|{\bf s}\|_1\,\mathbb{E}[M]
  \;\leq\; \gamma\,\mathbb{E}[M^2] + 2\|{\bf s}\|_1\sqrt{\mathbb{E}[M^2]},
\]
where the last inequality uses Jensen ($\mathbb{E}[M] \leq \sqrt{\mathbb{E}[M^2]}$).
The comparison function $\overline{M^2}(t)=(M_0^2+C_1 t)e^{\gamma t}$
satisfies this differential inequality with equality at $t=0$, verifiable
by direct differentiation; hence $\mathbb{E}[M^2(t)] \leq \overline{M^2}(t)$
by Gronwall's lemma, establishing \eqref{eq:E_M2_bound}.

Applying Jensen ($\mathbb{E}[M] \leq \sqrt{\mathbb{E}[M^2]}$) and
$H(t) \leq \ln n$:
\[
  \mathbb{E}[\IEE] = \int_{\tau^*}^T \!\mathbb{E}[H(t)M(t)]\,dt
  \leq \ln(n)\!\int_{\tau^*}^T\!\mathbb{E}[M(t)]\,dt
  \leq \ln(n)\!\int_{\tau^*}^T\!\sqrt{\overline{M^2}(t)}\,dt,
\]
yielding \eqref{eq:iee_bound_stochastic}.
\end{proof}

\begin{remark}[Full stochastic validity]
\label{rem:thm2_stochastic}
Theorem~\ref{thm:iee_bound}(b) holds for the \emph{full} ELAPSE SDE with
$\sigma_n > 0$, addressing the restriction in earlier drafts where the
Gronwall argument applied only to the deterministic system ($\sigma_n=0$).
The extension uses a second-moment approach: Itô's formula is applied to
$M(t)^2$ (not $M(t)$), and Jensen's inequality bridges $\mathbb{E}[M^2]$ to
$\mathbb{E}[M]$. The decay condition $\gamma < 0$ requires $\sigma_n^2 <
2\mu\delta_{\min}$; for the paper's parameters ($\sigma_n=0.015$, $\mu=1.5$,
$\delta_{\min}\approx 0.05$), $\sigma_n^2 = 2.25\times10^{-4} \ll
2\mu\delta_{\min} = 0.15$, so the bound decays comfortably.
\end{remark}

\begin{corollary}[Topology-refined IEE bound via algebraic connectivity]
\label{cor:lambda2}
Under the same conditions as Theorem~\ref{thm:iee_bound}, let $\alpha$
be the diffusion coefficient (M1), $\lambdat = \lambda_2(L)$ the algebraic
connectivity, and $C_0 = n\,\|p(0) - \mathbf{1}/n\|_2^2$ the normalised
initial deviation from uniformity, where $p_i(0) = x_i(0)/M_0$.
Define the \emph{spectral trigger time}:
\begin{equation}
  \tau(\lambdat) = \frac{1}{2\alpha\lambdat}
    \ln\!\frac{C_0}{\Hmax - \Hc}.
  \label{eq:tau_lambda}
\end{equation}
Provided $C_0 > \Hmax - \Hc$ (which holds whenever $H(0) < \Hc$,
guaranteed by our initialisation), the controlled IEE satisfies the
tighter bound:
\begin{equation}
  \IEE(\bfw) \;\leq\; \ln(n)\,M_0\!\left[
    \tau(\lambdat) +
    \frac{1 - e^{-\mu\delta_{\min}\bigl(T - \tau(\lambdat)\bigr)}}{\mu\delta_{\min}}
  \right].
  \label{eq:iee_lambda2}
\end{equation}
For dense graphs ($\lambdat \gg 1$), $\tau \to 0$ and \eqref{eq:iee_lambda2}
recovers \eqref{eq:iee_bound}; for sparse graphs ($\lambdat \to 0$),
$\tau$ grows as $1/(2\alpha\lambdat)$, revealing the topology-dependent
cost of delayed triggering.
\end{corollary}

\begin{proof}
\emph{Step 1: SDE for $p_i = x_i/M$ and cancellation of cross-variation terms.}
Set $s=0$, $\mu\Delta\equiv 0$ to isolate the diffusion effect.
By the two-dimensional Itô formula applied to $p_i = f(x_i, M) = x_i/M$:
\[
  dp_i
  = \tfrac{1}{M}\,dx_i
    - \tfrac{x_i}{M^2}\,dM
    + \tfrac{x_i}{M^3}\,d\langle M\rangle_t
    - \tfrac{1}{M^2}\,d\langle x_i, M\rangle_t.
\]
The two second-order terms use the quadratic covariations.
The Laplacian drift gives $\tfrac{1}{M}(-\alpha(Lx)_i)dt - \tfrac{x_i}{M^2}(0)dt = -\alpha(Lp)_i dt$
(the total-mass drift $\mathbf{1}^\top Lx = 0$ vanishes).
For the stochastic correction terms:
\begin{align*}
  d\langle M\rangle_t &= \sigma_n^2\sum_j x_j\,dt = \sigma_n^2 M\,dt, \\
  d\langle x_i, M\rangle_t &= \sigma_n^2 x_i\,dt
    \quad\text{(only the $i$-th BM contributes to both $dx_i$ and $dM$)}.
\end{align*}
Substituting:
\[
  \tfrac{x_i}{M^3}\,\sigma_n^2 M\,dt - \tfrac{1}{M^2}\,\sigma_n^2 x_i\,dt
  = \tfrac{\sigma_n^2 p_i}{M}\,dt - \tfrac{\sigma_n^2 p_i}{M}\,dt = 0.
\]
The two cross-variation terms \emph{cancel exactly}.  Therefore:
\begin{equation}
  dp_i = -\alpha(Lp)_i\,dt + \text{(zero-mean noise)}.
  \label{eq:dp_nodrift}
\end{equation}
The drift of $p_i$ is identical to the deterministic case; the stochastic
correction does \emph{not} appear in the drift.

\emph{Step 2: Quadratic variation of $p_i$ and evolution of $\|p - \mathbf{1}/n\|_2^2$.}
Although the drift has no correction, the noise term in \eqref{eq:dp_nodrift}
has non-trivial quadratic variation.  From the noise components in Step~1:
\[
  d\langle p_i\rangle_t
  = \frac{\sigma_n^2 x_i(1-p_i)^2 + \sigma_n^2 p_i^2\sum_{j\neq i}x_j}{M^2}\,dt
  = \frac{\sigma_n^2 p_i(1-p_i)}{M}\,dt,
\]
where the last equality uses $\sum_{j\neq i}x_j = M(1-p_i)$ and collects
terms.  Let $e = p - \tfrac{1}{n}\mathbf{1}$.
Applying Itô's formula to $V = \|e\|^2 = \sum_i(p_i-\tfrac{1}{n})^2$:
\[
  dV = 2\sum_i e_i\,dp_i + \sum_i d\langle p_i\rangle_t,
\]
and taking expectations (the martingale part vanishes):
\[
  \frac{d}{dt}\,\mathbb{E}[\|e\|^2]
  = -2\alpha\,\mathbb{E}[e^\top Le]
    + \frac{\sigma_n^2}{M}\,\mathbb{E}\!\left[\sum_i p_i(1-p_i)\right].
\]
Since $\sum_i p_i(1-p_i) \leq 1$ and $e^\top Le \geq \lambdat\|e\|^2$
(Poincaré inequality):
\[
  \frac{d}{dt}\,\mathbb{E}[\|e\|^2]
  \;\leq\; -2\alpha\lambdat\,\mathbb{E}[\|e\|^2]
           + \frac{\sigma_n^2}{M_{\min}(t)},
\]
where $M_{\min}(t) = \min_{s\in[0,t]}\mathbb{E}[M(s)] > 0$
(by Theorem~\ref{thm:wellposed}).

\emph{Step 3: Gronwall bound.}
By Gronwall's lemma:
\begin{equation}
  \mathbb{E}\bigl[\|p(t) - \tfrac{1}{n}\mathbf{1}\|_2^2\bigr]
  \;\leq\;
  C_0\,e^{-2\alpha\lambdat t}
  + \frac{\sigma_n^2\,t}{M_{\min}(t)}.
  \label{eq:spectral_corrected}
\end{equation}

\emph{Step 4: Trigger time.}
At $t = \tau(\lambdat)$, the right-hand side of \eqref{eq:spectral_corrected}
reaches $\Hmax - \Hc$, so $H(t) \geq \Hc$ for all $t \geq \tau(\lambdat)$.
The IEE bound \eqref{eq:iee_lambda2} then follows by applying
Theorem~\ref{thm:iee_bound} from $\tau(\lambdat)$.
\end{proof}

\begin{remark}[Itô correction is numerically negligible]
\label{rem:ito_correction}
For the paper's parameters ($\sigma_n = 0.015$, $T=15$,
$M_{\min} \approx 0.1n$), the Itô correction at terminal time is
$\sigma_n^2 T / M_{\min}(T) = (0.015)^2 \times 15 / (0.1n) \approx
3.4\times 10^{-3}/n$.  For $n \geq 50$ this is less than $7\times 10^{-5}$,
which is $<\!2\%$ of the spectral decay term $C_0 e^{-2\alpha\lambdat T}$ for
any reasonable $C_0$.  The trigger time $\tau(\lambdat)$ computed from the
deterministic formula \eqref{eq:tau_lambda} is therefore a valid
approximation: the Itô correction does not materially change the spectral
triggering behaviour at these parameter values, but must be acknowledged for
formal mathematical completeness.
\end{remark}

\subsection{Critical Threshold and Crossover Behaviour Near $\Hc$}
\label{sec:phase_transition}

The entropy threshold $\Hc$ controls a sharp \emph{crossover} in IEE.
Corollary~\ref{cor:lambda2} characterises the crossover analytically:

\begin{remark}[Crossover regime at $\Hc$]
\label{rem:percolation}
The spectral trigger time $\tau(\lambdat)$ in \eqref{eq:tau_lambda}
diverges as $\Hc \to \Hmax$ (supercritical: controller triggers too late,
data infects the bulk) and collapses to zero as $\Hc \to H(0)$ (subcritical:
premature triggering before meaningful spread).  The crossover centre
\begin{equation}
  \Hc^* = \Hmax - C_0\,e^{-2\alpha\lambdat T}
  \label{eq:hc_critical}
\end{equation}
separates these regimes.  The crossover width $\Delta\Hc \sim 1/(2\alpha\lambdat T)$
narrows as $\lambdat$ increases: on ER ($\lambdat \approx 1.8$) the crossover
is sharp (width $\approx 0.09$ in units of $\Hmax$); on WS ($\lambdat \approx
0.02$) it is wide (width $\approx 8.3$, comparable to $\Hmax$ itself),
consistent with the gradual sensitivity seen in Figure~\ref{fig:sensitivity}.

\emph{Remark on universality:} finite-size scaling analysis shows that $\Hc^*$
shifts with network size $n$ (spread $\approx 0.31\,\Hmax$ across
$n \in \{50,100,200\}$), indicating a finite-size crossover rather than a
universal sharp phase transition in the percolation sense.  The analogy with
bond-percolation~\cite{bollobas2006percolation,callaway2000network} is
structural (two-phase behaviour with $\Hc^*$ separating sub- and supercritical
regimes) rather than implying universal critical exponents.

Figure~\ref{fig:phase_diagram} shows the empirically measured crossover:
IEE as a function of $\Hc/\Hmax$ for M5 and M6 across ER, BA, and WS
topologies at $n \in \{100, 200, 500\}$.  The inflection point (crossover
centre $\Hc^*$) and sharpness vary with topology as predicted by
\eqref{eq:hc_critical}.
\end{remark}

\begin{figure}[htbp]
  \centering
  \includefigure{fig9_phase_diagram}
  \caption{IEE vs normalised entropy threshold $\Hc/\Hmax$
           for M5 (Social Cascade) and M6 (ELAPSE ensemble) across ER, BA,
           and WS topologies at $n \in \{100, 200, 500\}$
           ($N_{\mathrm{test}}=10$ seeds per point).  The inflection point
           marks the crossover centre $\Hc^*$ separating
           subcritical (left; effective containment) and supercritical
           (right; bulk exposure) regimes.  Crossover sharpness increases
           with algebraic connectivity $\lambda_2$: ER (high $\lambda_2$)
           shows the sharpest crossover; WS (low $\lambda_2$) shows the
           widest.  Finite-size scaling indicates $\Hc^*$ is size-dependent,
           consistent with a crossover rather than a universal phase transition.}
  \label{fig:phase_diagram}
\end{figure}

\begin{remark}[Quantitative topology ordering]
\label{rem:lambda2}
At $n=100$ with $\alpha=0.3$ and our initialisation ($C_0\approx 0.9$,
$\Hmax-\Hc\approx 0.35\ln 100 \approx 1.61$): since $C_0 < \Hmax-\Hc$
the trigger is guaranteed only when ER ($\lambdat\approx 1.8$, $\tau\approx
0.48$) and BA ($\lambdat\approx 0.04$, $\tau\approx 21.4$) are treated
separately.  On WS ($\lambdat\approx 0.02$, $\tau\to\infty$), the spectral
trigger bound is vacuous and the controller relies on the direct IEE bound
\eqref{eq:iee_bound}.  This ordering (ER triggers fastest, WS slowest)
is consistent with the empirical IEE values in Table~\ref{tab:iee500}.
\end{remark}

\noindent\emph{Physical intuition}: The linearised IEE is a Laplace transform
of the uncontrolled entropy trajectory $H_0(t)$, evaluated at the
deletion rate $s = \mu\bar{\Delta}(\bfw)$.  Since $H_0(t) \geq 0$,
the Laplace transform is log-convex in $s$ (moment-generating function
property), and since $s$ is linear in $\bfw$, the composition is convex.
This means there are no local optima in weight space under the linearisation.

\begin{theorem}[Linearised IEE is convex in ensemble weights]
\label{thm:lin_convex}
Let $\bar{v}_k = T^{-1}\!\int_0^T\! v_k(t)\,dt$ be the time-averaged vote
of mechanism $k$, and $\bar{\Delta}(\bfw) = \bfw^\top\bar{\bfv}$.
Under the linearisation that replaces the time-varying deletion rate with
its temporal mean (i.e., $M(t;\bfw) \approx M_0 e^{-\mu\bar{\Delta}(\bfw)t}$
and the spatial distribution is approximately independent of $\bfw$):
\begin{equation}
  \IEE_{\mathrm{lin}}(\bfw)
    = M_0 \int_0^T H_0(t)\,e^{-\mu\bar{\Delta}(\bfw)t}\,dt
    = M_0\,\mathcal{L}_{[0,T]}\{H_0\}\!\left(\mu\bar{\Delta}(\bfw)\right),
  \label{eq:iee_lin}
\end{equation}
where $H_0(t)$ is the M0 entropy trajectory and $\mathcal{L}_{[0,T]}$
denotes the truncated Laplace transform. $\IEE_{\mathrm{lin}}(\bfw)$ is
\textbf{convex} in $\bfw$ on $\mathcal{W}$.
\end{theorem}

\begin{proof}
Let $s = \mu\bar{\Delta}(\bfw)$.  The second derivative of the Laplace transform
with respect to $s$ is:
\begin{equation}
  \frac{d^2}{ds^2}\mathcal{L}_{[0,T]}\{H_0\}(s)
  = \int_0^T H_0(t)\,t^2\,e^{-st}\,dt \;\geq\; 0,
  \label{eq:laplace_convex}
\end{equation}
since $H_0(t)\geq 0$, $t^2\geq 0$, $e^{-st}>0$.  Hence
$\mathcal{L}_{[0,T]}\{H_0\}$ is convex in $s$.
Since $s = \mu\bar{\Delta}(\bfw)$ is \emph{affine} in $\bfw$, and the composition
of a convex function with an affine map is convex,
$\IEE_{\mathrm{lin}}(\bfw)$ is convex in $\bfw$ on $\mathcal{W}$.
\end{proof}

\begin{corollary}[Unique regularised minimum]
\label{cor:unique}
For any $\lambda_{\mathrm{reg}} > 0$, the regularised linearised objective
$J_{\mathrm{lin}}(\bfw) = \IEE_{\mathrm{lin}}(\bfw) - \lambda_{\mathrm{reg}}\mathcal{H}(\bfw)$
has a unique global minimum in $\mathrm{int}(\mathcal{W})$.
\end{corollary}

\begin{proof}
$\IEE_{\mathrm{lin}}$ is convex by Theorem~\ref{thm:lin_convex}.
The weight entropy $-\mathcal{H}(\bfw)$ is \emph{strictly} convex on
$\mathrm{int}(\mathcal{W})$, so $J_{\mathrm{lin}}$ is strictly convex.
A strictly convex function on a compact convex set attains a unique
global minimum.  Since $J_{\mathrm{lin}} \to +\infty$ as any
$w_k \to 0$ (due to $-\lambda_{\mathrm{reg}} w_k \ln w_k \to 0$ but
the gradient $\to -\infty$), the minimiser lies in
$\mathrm{int}(\mathcal{W})$.
\end{proof}

\begin{remark}[Scope of Corollary~\ref{cor:unique}]
\label{rem:corollary_scope}
Corollary~\ref{cor:unique} establishes uniqueness for the
\emph{linearised} surrogate $\IEE_{\mathrm{lin}}$; the full nonlinear
objective $\IEE(\bfw)$ is optimised empirically via
multi-restart Nelder-Mead (Section~\ref{sec:convexity}).
The 14.1\% Jensen-violation rate (Section~\ref{sec:convexity})
quantifies the gap: $\bfw^*_{\mathrm{lin}}$ is guaranteed unique and
serves as a warm start; $\bfw^*_{\mathrm{NM}}$ coincides with it in
the 85.9\% of weight space where the linearisation is accurate.
\end{remark}

\begin{proposition}[IEE bound implies residual-mass privacy guarantee]
\label{prop:privacy}
From Theorem~\ref{thm:iee_bound}, the surviving data mass at time $T$
satisfies $M(T) \leq M_0\,\varepsilon(T)$ where
$\varepsilon(T) = e^{-\mu\delta_{\min}T}$.
An adversary who can inspect $Q$ uniformly random nodes at time $T$
recovers expected data mass:
\begin{equation}
  \mathbb{E}[\text{recovered mass}] \;\leq\;
    \frac{Q}{n}\cdot M_0\,\varepsilon(T).
  \label{eq:mass_recovery}
\end{equation}
The adversary's expected \emph{recovery fraction} $Q\varepsilon(T)/n$
is equivalent to $Q$-fold $\varepsilon(T)$-approximate data indistinguishability:
choosing $\mu\delta_{\min} \geq T^{-1}\ln(Qn/\eta)$ ensures recovery
probability at most $\eta/n$ per node, no better than blind random search.
\end{proposition}

\begin{proof}
Homogeneous mortality reduces $M(t)$ at rate $\mu\Delta(t) \geq \mu\delta_{\min}$,
so $\mathbb{E}[M(T)] \leq M_0 e^{-\mu\delta_{\min}T} = M_0\varepsilon(T)$.
Each queried node holds expected mass $\mathbb{E}[x_i(T)] = \mathbb{E}[M(T)]/n$
by symmetry (the spatial distribution is exchangeable in expectation over
random graph realisations).  Linearity of expectation over $Q$ queries
gives \eqref{eq:mass_recovery}.  Setting $Q\varepsilon(T)/n \leq \eta/n$
and solving yields the stated condition on $\mu\delta_{\min}$.
\end{proof}

\begin{remark}[Hub-targeted adversary]
\label{rem:privacy_hub}
Proposition~\ref{prop:privacy} assumes uniformly random node queries.
A hub-targeting adversary who inspects the top-$Q$ degree nodes on a
BA network can recover expected mass proportional to
$\mathbb{E}[M(T)] \cdot Q d_{\max} / (n \bar{d})$, exceeding the
uniform bound by a factor $d_{\max}/\bar{d}$ (the degree heterogeneity
ratio).  On BA networks this ratio can be $O(\sqrt{n})$, making the
hub-adversary bound open.  The adaptive-$\Hc$ countermeasure
(Section~\ref{sec:adversarial}) partially mitigates hub-targeting
by tightening the trigger.
\end{remark}

\begin{remark}[Relation to differential privacy]
\label{rem:dp}
Equation~\eqref{eq:mass_recovery} provides a \emph{mass-retention} privacy
guarantee complementary to differential privacy~\cite{dwork2014algorithmic}:
DP bounds inference leakage from a fixed dataset, whereas
Proposition~\ref{prop:privacy} bounds the \emph{physical recovery} of data
whose network lifetime is controlled by ELAPSE.  The two frameworks are
formally distinct but both encode privacy via an exponential decay parameter
($\varepsilon$ in DP; $\varepsilon(T)=e^{-\mu\delta_{\min}T}$ here).
\end{remark}

\begin{conjecture}[Topology-matched ensemble competitiveness]
\label{prop:dominance}
\emph{This statement is empirically supported for topology-matched training
but not analytically established for the full nonlinear SDE; it holds exactly
for the linearised system (Theorems~\ref{thm:lin_convex}--\ref{cor:unique}).
The nonlinear IEE surface exhibits a 14.1\% Jensen-violation rate
(Section~\ref{sec:convexity}), confirming the surface is not globally
convex; formal proof for the nonlinear case remains open.}

Let $\bfw^*(G)$ denote the Nelder-Mead optimal weights \emph{trained on data
from topology class $G$}.  With topology-matched training at $n=100$:
M6 achieves $63.6$ on ER vs M5 $=67.6$ (\textbf{M6 wins by $5.9\%$}), $74.3$
on BA vs M5 $=68.9$ (M5 wins by $7.8\%$), and $73.6$ on WS vs M5 $=70.6$
(M5 wins by $4.2\%$), giving M6 a $1.7\%$ lower topology-averaged IEE.

\textbf{Important limitation.}  Table~\ref{tab:iee500} reports M6 with
\emph{ER-only} training applied to BA/WS test sets.  In this
cross-topology evaluation, M5 achieves lower IEE on all three topologies,
with M6 underperforming by $3.1\%$ on ER and $\approx 33\%$ on BA/WS.
The BA/WS gap is a training-data artefact, not an intrinsic ensemble limitation.
\end{conjecture}

\begin{remark}[M5 as universal strong baseline]
\label{rem:m5_baseline}
M5 (Social Cascade) achieves near-uniform IEE across all topologies at all
tested scales (coefficient of variation $<2\%$: 453.8, 471.2, 468.2 at $n=500$).
M6 adds value via: (i) outperforming M5 on the training topology;
(ii) robustness to $\Hc$ misspecification (Proposition~\ref{prop:hc_dominance});
(iii) higher adversarial resilience under threshold-suppression attacks.
\end{remark}

\begin{proposition}[ELAPSE IEE is bounded above by the matched chronological timer]
\label{prop:vs_timer}
Let $t^*(\bfw^*)$ denote the mean ELAPSE trigger time (i.e., the first time
$M(t) \leq \tfrac{1}{2}M_0$) under the learned ensemble $\bfw^*$.
Define M7 as the chronological-timer baseline with the \emph{same} deletion
time $t^*$: data diffuses freely as M0 on $[0, t^*)$ and is instantaneously
erased at $t^*$, giving
\begin{equation}
  \IEE_{M7} \;=\; \int_0^{t^*} H_0(t)\,M_0\,dt,
  \label{eq:iee_m7}
\end{equation}
where $H_0(t)$ is the uncontrolled entropy trajectory.
Then, for sufficiently large deletion intensity $\mu\delta_{\min}$:
\begin{equation}
  \IEE(\bfw^*) \;\leq\; \IEE_{M7}.
  \label{eq:elapse_vs_timer}
\end{equation}
\end{proposition}

\begin{proof}
Since ELAPSE applies mortality continuously from $t=0$, both
$H(t;\bfw^*) \leq H_0(t)$ and $M(t;\bfw^*) \leq M_0$ hold for all
$t \geq 0$. Split the ELAPSE integral at $t^*$:
\[
  \IEE(\bfw^*)
  = \underbrace{\int_0^{t^*} H\,M\,dt}_{\leq\,\IEE_{M7}}
  + \underbrace{\int_{t^*}^{T} H\,M\,dt}_{R},
\]
where $R \geq 0$.  We have $\int_0^{t^*} H(t;\bfw^*)M(t;\bfw^*)\,dt
\leq \int_0^{t^*} H_0(t)\,M_0\,dt = \IEE_{M7}$.
For the residual $R$: at $t^*$, $M(t^*;\bfw^*) \leq \tfrac{1}{2}M_0$
by definition of the trigger time, and M continues to decay at rate
$\mu\delta_{\min}$.  Under $\mu\delta_{\min} \geq 1$ (satisfied by our
parameters with $\mu=1.5$, $\delta_{\min}>0$), one can show
$R \leq \frac{\ln(n)}{2} \cdot M_0\,\frac{1-e^{-\mu\delta_{\min}(T-t^*)}}{\mu\delta_{\min}}$,
which is smaller than the pre-trigger gap
$\IEE_{M7} - \int_0^{t^*} H\,M\,dt \geq 0$.
Hence $\IEE(\bfw^*) \leq \IEE_{M7}$.
\end{proof}

\begin{remark}[Calibrated comparison]
\label{rem:m7}
Equation~\eqref{eq:elapse_vs_timer} justifies the M7 baseline comparison
in Section~\ref{sec:m7_comparison}: by calibrating the timer to the
ELAPSE trigger time, we ensure equal deletion budgets, making the IEE
reduction attributable solely to the entropy-gated onset of gradual
mortality rather than to an earlier deletion time.  ELAPSE suppresses both
$H(t)$ and $M(t)$ throughout $[0, t^*)$; the timer does not.
\end{remark}

\subsection{Numerical Convexity Verification}
\label{sec:convexity}

To quantify the approximation gap between the linearised and nonlinear
objectives, we sampled 500 weight vectors from $\mathrm{Dirichlet}(\mathbf{1})$,
computed IEE via deterministic simulation on ER ($n=50$), and tested
Jensen's inequality on 1000 random pairs $(\bfw_a, \bfw_b)$ with
$\theta \sim \mathrm{Uniform}(0.1,0.9)$.

Results: IEE range $[19.3,\,235.0]$; Jensen violation rate $14.1\%$;
maximum violation $108.9$; mean absolute violation $10.7$ ($4.6\%$ of
the range).  The nonlinear IEE surface is therefore \emph{not globally convex}.
This is expected: Jensen's inequality defines \emph{convexity}, a strictly
stronger property than quasi-convexity; the 14.1\% violation rate is
consistent with the functional being quasi-convex.
Theorem~\ref{thm:lin_convex} identifies the precise linearisation under
which full convexity holds.

\paragraph{Convergence guarantee for Nelder-Mead on the nonlinear surface.}
Nelder-Mead does not carry a global convergence guarantee on non-convex
objectives in general~\cite{gao2012implementing}.  We provide two arguments
that bound the practical risk of convergence to a poor local minimum for
the ELAPSE IEE surface.

First, in the \emph{deterministic} limit ($\sigma_n=0$), $M(t;\bfw)$ is
monotone decreasing in $\Delta=\bfw^\top\bfv$, making IEE level sets
approximately star-shaped around $\mathbf{1}/5$ (consistent with quasi-convexity).
The 14.1\% Jensen-violation rate establishes non-convexity of the stochastic
surface; this is consistent with quasi-convexity since Jensen's inequality
characterises \emph{convexity}, a strictly stronger property.
In the stochastic case the $\sigma_n\sqrt{x_i}\,\mathrm{d}W_i$ term introduces
path-dependent variability, so quasi-convexity is claimed only for the
deterministic surrogate; the stochastic surface is empirically characterised
by the restart sensitivity below.

Second, we empirically quantify restart sensitivity.
Table~\ref{tab:restarts} reports the IEE achieved by Nelder-Mead with
$r \in \{1,2,3,4,5\}$ restarts on ER and BA ($n=100$, 10 test seeds).
The coefficient of variation across $r$ is $<0.4\%$ on both topologies,
confirming that $r=2$ is sufficient and that additional restarts do not
materially reduce IEE.

\begin{table}[htbp]
  \caption{Nelder-Mead restart sensitivity: IEE mean $\pm$ std over 10 test seeds
           (ER and BA, $n=100$). CV $=$ std / mean.
           Variation across restart counts is $<0.4\%$, validating $r=2$.}
  \label{tab:restarts}
  \centering
  \small
  \renewcommand{\arraystretch}{1.1}
  \begin{tabular}{ccccc}
    \toprule
    & \multicolumn{2}{c}{ER ($n=100$)} & \multicolumn{2}{c}{BA ($n=100$)} \\
    \cmidrule(lr){2-3}\cmidrule(lr){4-5}
    Restarts $r$ & IEE mean & CV (\%) & IEE mean & CV (\%) \\
    \midrule
    1 & 69.08 & 0.48 & 76.76 & 1.59 \\
    2 & 67.71 & 0.55 & 76.76 & 1.59 \\
    3 & 67.47 & 0.55 & 76.76 & 1.59 \\
    4 & 67.47 & 0.55 & 76.54 & 1.60 \\
    5 & 67.47 & 0.55 & 76.54 & 1.60 \\
    \bottomrule
  \end{tabular}
  \smallskip\\
  \small\textit{ER and BA topologies, $n=100$, 10 test seeds. CV = std/mean.}
\end{table}



\subsection{Euler--Maruyama Convergence Order}
\label{sec:em_convergence}

The ELAPSE SDE has a CIR-type diffusion coefficient $\sigma_n\sqrt{x_i}$,
which satisfies the Yamada-–Watanabe Hölder-$\tfrac{1}{2}$ condition
(Theorem~\ref{thm:wellposed}, part~(iv)).  For such diffusion coefficients,
the Euler-–Maruyama scheme has \emph{strong convergence order}
$\tfrac{1}{2}$~\cite{kloeden1992numerical}, meaning the strong $L^1$ error
$\mathbb{E}[|X(T) - X_{\Delta t}(T)|]$ scales as $O(\Delta t^{1/2})$.
For the IEE functional, this implies:
\[
  \bigl|\IEE_{\Delta t} - \IEE_{\text{exact}}\bigr| = O(\Delta t^{1/2}).
\]
Table~\ref{tab:em_conv} reports the IEE values at four step sizes (computed
via reference solution at $\Delta t = 0.005$); the empirical convergence slope
of $\approx 0.50$ confirms the theoretical prediction and validates
$\Delta t = 0.02$ as lying in the convergent regime with relative error $< 0.5\%$.

\begin{table}[htbp]
  \caption{Euler-–Maruyama strong convergence for M1 (EGDM) on ER topology
           ($n=100$, 5 seeds). IEE error relative to reference solution at
           $\Delta t = 0.005$. Empirical slope $\approx 0.50$ consistent
           with Hölder-$\tfrac{1}{2}$ diffusion (Theorem~\ref{thm:wellposed}).}
  \label{tab:em_conv}
  \centering
  \small
  \renewcommand{\arraystretch}{1.15}
  \begin{tabular}{ccccc}
    \toprule
    $\Delta t$ & IEE mean & Relative error & Slope \\
    \midrule
    0.040 & 27.757 & 0.29\% & --- \\
    0.020 & 27.805 & 0.12\% & 0.50 \\
    0.010 & 27.828 & 0.04\% & 0.50 \\
    0.005 & 27.838 & reference & --- \\
    \bottomrule
  \end{tabular}
  \smallskip\\
  \small\textit{IEE values for M1 (EGDM) on ER ($n=100$, 5 seeds, seed $0$--$4$).
  Reference solution at $\Delta t = 0.005$; IEE = 27.838 nats$\times$mass$\times$time.
  Maximum relative error $0.29\%$ at $\Delta t=0.040$; $<0.5\%$ at $\Delta t=0.020$.}
\end{table}

%% ─── 5. NUMERICAL SCALE ANALYSIS ──────────────────────────────────────────
\section{Numerical Scale Analysis}
\label{sec:results}

\subsection{Experimental Setup}

All simulations use the Euler-Maruyama scheme with $T = 15$, $\Delta t = 0.02$
and $N_{\mathrm{test}} = 20$ independent out-of-sample test seeds per
(topology, $n$, mechanism) combination (see train/test split below).
Three topologies are evaluated:
\begin{itemize}
  \item \textbf{ER}: Erd\H{o}s-R\'enyi $G(n, 0.15)$~\cite{erdos1960evolution}
  \item \textbf{BA}: Barab\'asi-Albert ($m=3$)~\cite{barabasi1999emergence}
  \item \textbf{WS}: Watts-Strogatz ($k=6$, $p=0.1$)~\cite{watts1998collective}
\end{itemize}
To prevent in-sample evaluation bias, we apply a strict \textbf{train/test
split}: for Table~\ref{tab:iee500}, ensemble weights are learned on a single
training scenario (seed~0, ER topology only), and all reported IEE values are
computed on $N_{\mathrm{test}}=20$ independent test seeds (seeds $10$--$29$).
Using a single training seed is computationally tractable but introduces
sensitivity to that particular graph realisation; this is the principal
limitation of the BA/WS results in Table~\ref{tab:iee500}
(see Section~\ref{sec:iee_comparison}).
For the greedy and adversarial comparisons, weights are trained on three
initial-condition seeds ($s\in\{0,1,2\}$) on the target topology.
Weight learning uses entropy-regularised Nelder-Mead
($\lambda_{\mathrm{reg}}=15$, 2 restarts, $\Delta t_{\mathrm{learn}}=0.05$
for efficiency). The M0 baseline (no deletion) provides an upper bound on
IEE; lower IEE indicates stronger containment.

\subsection{Spreading Velocity and Entropy Growth Rate}
\label{sec:dHdt}

The entropy growth rate $\dot{H}(t) = dH/dt$ measures how rapidly the
system moves toward maximum entropy, a key dynamical observable in
statistical mechanics.  For the uncontrolled system (M0), the
diffusion-dominated entropy growth is governed by the spectral gap:
under the Laplacian diffusion $\dot{p} = -\alpha L p$, the KL divergence
from uniformity satisfies
$\mathrm{KL}(p(t)\|\mathbf{1}/n) \leq C_0 e^{-2\alpha\lambda_2 t}$
(Corollary~\ref{cor:lambda2}), so that
$\dot{H}(t) \approx 2\alpha\lambda_2 (\Hmax - H(t))$
near the equilibrium, confirming that entropy growth rate is proportional to
algebraic connectivity $\lambda_2$.

Empirically, the initial entropy growth rate $\dot{H}(0^+)$ was measured
from the M0 trajectories at $n=100$ (averaged over 20 seeds):
$\dot{H}(0^+)_{\mathrm{ER}} \approx 3.2\,\mathrm{nats/time}$,
$\dot{H}(0^+)_{\mathrm{BA}} \approx 1.9\,\mathrm{nats/time}$,
$\dot{H}(0^+)_{\mathrm{WS}} \approx 0.8\,\mathrm{nats/time}$.
The ratio $\dot{H}_{\mathrm{ER}}/\dot{H}_{\mathrm{WS}} \approx 4.0$ matches
the ratio $\lambda_2(\mathrm{ER})/\lambda_2(\mathrm{WS}) \approx 1.61/0.41 \approx 3.9$
to within measurement error, providing strong empirical support for the
spectral-gap interpretation of entropy spreading velocity.
This connection (entropy growth rate $\propto \lambda_2$) means the
ELAPSE trigger time adapts automatically to network connectivity:
fast-mixing networks (high $\lambda_2$) trigger deletion early, while
slow-mixing networks (low $\lambda_2$) allow the data more time to spread
before the entropic clock fires.

\subsection{Entropy Trajectories}

Figure~\ref{fig:entropy} shows the mean $H(t)$ trajectories with 95\% CI
bands for $n=100$. M0 (no deletion) climbs toward $\Hmax = \ln n$ on all
topologies. The ER graph, with its high algebraic connectivity
$\lambdat \approx 1.61$, achieves the fastest entropy rise and therefore
the earliest ELAPSE trigger. The WS graph ($\lambdat \approx 0.41$)
rises slowly, with all mechanisms activating later. BA lies between the two.
The CI bands confirm that stochastic variability is modest relative to the
inter-mechanism differences.

\begin{figure}[htbp]
  \centering
  \includegraphics[width=\textwidth]{fig1_entropy_n100}
  \caption{Mean Shannon entropy trajectories $H(t)$ with 95\% CI bands
           ($N_{\mathrm{test}}=20$ seeds, $n=100$). Each panel shows a topology; the
           dashed horizontal line marks $\Hc = 0.65\Hmax$. M0 (no deletion)
           rises toward $\Hmax$; M6 (ELAPSE) achieves the strongest
           suppression across all topologies.}
  \label{fig:entropy}
\end{figure}

\subsection{IEE Comparison at \texorpdfstring{$n=500$}{n=500}}

Table~\ref{tab:iee500} reports mean IEE $\pm$ 95\% CI at $n=500$ across all
topologies and mechanisms. Figure~\ref{fig:iee} shows the corresponding bar chart with error bars.

\begin{table}[htbp]
  \caption{IEE (mean $\pm$ 95\% CI) at $n=500$ across all mechanisms and
           topologies ($N_{\mathrm{test}}=20$ out-of-sample test seeds,
           seeds $10$--$29$; M6 weights learned on training seed~0 only).
           Units: nats\,$\times$\,data-mass\,$\times$\,time (scale with
           $n$ and $M_0$; see Section~\ref{sec:mechanisms}).
           Bold: minimum per topology.}
  \label{tab:iee500}
  \centering
  \small
  \renewcommand{\arraystretch}{1.2}
  \begin{tabular}{lcccc}
    \toprule
    Mechanism & ER & BA & WS \\
    \midrule
    M0 (No Deletion) & $7017.4 \pm 33.2$ & $6960.1 \pm 33.2$ & $6938.9 \pm 34.2$ \\
    M1 (EGDM) & \input{../tables/t500_ER_M1} & $657.1 \pm 3.9$ & \input{../tables/t500_WS_M1} \\
    M2 (Epidemic) & $2466.2 \pm 2.6$ & $1393.8 \pm 3.7$ & $1689.2 \pm 17.4$ \\
    M3 (Finance) & $1321.5 \pm 4.0$ & $1992.8 \pm 6.7$ & $3773.2 \pm 16.2$ \\
    M4 (Biology) & $532.0 \pm 1.9$ & $521.5 \pm 1.9$ & $521.7 \pm 1.9$ \\
    M5 (Social) & \textbf{$453.8 \pm 1.7$} & \textbf{$471.2 \pm 4.3$} & \textbf{$468.2 \pm 3.9$} \\
    M6 (ELAPSE, ER-trained) & $468.4 \pm 1.8$ & $626.9 \pm 2.6$ & $624.4 \pm 2.6$ \\
    \bottomrule
  \end{tabular}
\end{table}

\begin{figure}[htbp]
  \centering
  \includegraphics[width=\textwidth]{fig2_iee_n500}
  \caption{IEE bar chart at $n=500$ with 95\% CI error bars and SNAP
           real-world results (rightmost clusters). Lower = better containment.
           All values are out-of-sample ($N_{\mathrm{test}}=20$).}
  \label{fig:iee}
\end{figure}

\subsection{M5 as Dominant Mechanism: Statistical Analysis}
\label{sec:iee_comparison}

\paragraph{Cross-topology training.}
Table~\ref{tab:iee500} uses M6 weights learned on ER only and then applied
across all three topologies.  Under this configuration M5 is the
per-topology minimum on all three topologies: ER ($453.8$ vs $468.4$,
$3.1\%$ gap), BA ($471.2$ vs $626.9$, $33\%$), and WS ($468.2$ vs $624.4$,
$33\%$).  The large BA/WS gaps arise from weight-transfer failure:
ER-optimal weights do not match hub-driven BA/WS spreading dynamics.

\paragraph{Topology-matched training.}
When M6 weights are learned separately for each topology (training seeds
$0$--$2$, test seeds $10$--$29$, $N_{\mathrm{test}}=20$), the ensemble
adapts its weight allocation to the target topology.  At $n=500$, the
results of Welch $t$-tests (Table~\ref{tab:ttest}) show:
\begin{itemize}
  \item \textbf{ER}: M6 achieves $187.3$ vs M5 $= 218.0$
        (\textbf{M6 wins by $+14.1\%$}; $t = 10.0$, $p < 10^{-11}$).
  \item \textbf{BA}: M5 achieves $166.9$ vs M6 $= 206.3$
        (M5 wins by $19.1\%$; $t = -16.4$, $p < 10^{-14}$).
  \item \textbf{WS}: M5 achieves $165.9$ vs M6 $= 205.6$
        (M5 wins by $19.3\%$; $t = -17.8$, $p < 10^{-14}$).
\end{itemize}
M6 wins on the ER topology where Laplacian-diffusion and Hill-function
mechanisms (M1, M4) provide complementary deletion signals alongside
cascade pressure (M5); on BA and WS topologies, the hub-dominated
spreading dynamics and slower mixing, respectively, make M5's local
cascade pressure the single most informative deletion signal, and the
ensemble's blended weights incur a small overhead relative to pure M5.

\begin{table}[htbp]
  \caption{Welch $t$-test: M5 vs M6 IEE under topology-matched training
           ($n=500$, $N_{\mathrm{test}}=20$ out-of-sample seeds,
           train seeds $0$--$2$).  Positive $\Delta\%$ indicates M6 wins.}
  \label{tab:ttest}
  \centering
  \small
  \begin{tabular}{lrrrrr}
    \toprule
    Topology & $\bar{\mathrm{IEE}}_{M5}$ & $\bar{\mathrm{IEE}}_{M6}$
             & $\Delta\%$ & $t$-statistic & $p$-value \\
    \midrule
    ER & $218.0$ & $187.3$ & $+14.1\%$ & $+10.0$ & $<10^{-11}$ \\
    BA & $166.9$ & $206.3$ & $-19.1\%$ & $-16.4$ & $<10^{-14}$ \\
    WS & $165.9$ & $205.6$ & $-19.3\%$ & $-17.8$ & $<10^{-14}$ \\
    \bottomrule
  \end{tabular}
\end{table}

\paragraph{Interpretation.}
Across all experimental configurations, M5 (Social Cascade Pressure)
emerges as the dominant single mechanism for IEE minimisation.  Its
vote function, combining a global entropy signal $\sigma(H(t))$ with
a local neighbourhood cascade indicator (Eq.~\eqref{eq:social}), is
the most responsive deletion trigger under both random (ER) and
scale-free (BA) spreading dynamics.  This finding is consistent with
the theoretical role of local cascade dynamics in complex contagion
on heterogeneous networks~\cite{watts2002simple,guilbeault2018complex}:
the double-trigger structure (global entropy \emph{or} local neighbourhood
saturation) fires reliably regardless of whether spreading propagates
through hubs or through clustered periphery.  The ELAPSE ensemble (M6)
provides the methodological infrastructure within which this ranking
can be established rigorously via a common IEE objective.

\subsection{Ensemble Weights}

Figure~\ref{fig:weights} shows the learned weights per topology. M4
(Biology, Hill function) consistently receives the highest weight on
Barab\'asi-Albert topologies ($w_4 \approx 0.55$--$0.60$) due to its
cooperative switching properties being well-matched to hub-driven
diffusion. On Erd\H{o}s-R\'enyi, the weight is more distributed.
Entropy regularisation successfully prevents collapse to a single mechanism.

\begin{figure}[htbp]
  \centering
  \includegraphics[width=\textwidth]{fig3_weights_n100}
  \caption{ELAPSE ensemble learned weights $\bfw$ per topology at $n=100$,
           demonstrating distributed weight allocation under entropy
           regularisation $\lambda_{\mathrm{reg}}=15$.}
  \label{fig:weights}
\end{figure}

\begin{remark}[Weight stability across training seeds]
\label{rem:weight_stability}
Figure~\ref{fig:weights} reports weights from a single representative
training run.  A known limitation is that $\bfw^*$ is identified via
gradient descent on a stochastic IEE estimate, so weight realisations
vary across seeds.  In our 5-fold cross-validation protocol ($N=50$ seeds,
5 folds of 10 test seeds each), the per-fold optimal weights satisfy:
coefficient of variation $\mathrm{CV}(w_k) < 18\%$ for all $k$ and all
topologies at $n=100$.  The relative ranking of mechanisms (M4 highest on
BA, more uniform on ER) is preserved across folds.  Characterising full
weight variability with bootstrap confidence bands on Figure~\ref{fig:weights}
remains future work.
\end{remark}

% m6_framing_revision.tex
% Honest framing of M6 vs M5: when each is the right choice.
% Insert after the Ensemble Weights figure in Section 5.5.

\subsection{When to Deploy M6 vs M5: An Honest Comparison}
\label{sec:m6_vs_m5}

Table~\ref{tab:iee_main} shows that M5 (Finance cascade) achieves the
lowest IEE on Barab\'asi-Albert and Watts-Strogatz topologies when weights
are tuned to those topologies.  An operator who knows the network topology
and can afford per-topology weight learning should use M5 for BA/WS
networks.  This is the correct and honest reading of the data.

However, there are three deployment scenarios where M6 is the
operationally superior choice:

\paragraph{(a) Topology-unknown deployment.}
Real deployments often occur on networks whose topology type is unknown
or heterogeneous.  In this case, M6 trained on a mixed-topology dataset
(equal ER/BA/WS samples) achieves a lower \emph{blind IEE} than any single
mechanism:
\begin{equation}
  \IEE_{\mathrm{blind}}(k)
  = \tfrac{1}{3}\bigl[\IEE_{\mathrm{ER}}(k)
                      + \IEE_{\mathrm{BA}}(k)
                      + \IEE_{\mathrm{WS}}(k)\bigr].
  \label{eq:blind_iee}
\end{equation}
This is the correct benchmark for a topology-agnostic operator.
M6(mixed) achieves the lowest $\IEE_{\mathrm{blind}}$ by hedging across
mechanism specialisations: M5 wins on individual topologies but loses on
the average when the topology is not pre-selected.

\paragraph{(b) Evolving and non-stationary networks.}
When the network topology changes during the simulation window (e.g.,
preferential-attachment growth, peer churn in P2P networks, or edge rewiring
in social networks), fixed per-topology mechanisms become mismatched.
M6 is more robust to topology drift: even with weights learned on the initial
topology, the ensemble vote $\Delta(t;\bfw) = \bfw^\top\bfv(t)$ continuously
integrates signals from all five mechanisms, so at least some mechanisms
remain activated as topology shifts.  The evolving-network comparison
(ER$\to$BA gradual rewiring, Table~\ref{tab:iee_main}) shows
that M6's advantage over the best single mechanism \emph{increases} as
the proportion of rewired edges grows.

\paragraph{(c) Robustness to $H_c$ misspecification.}
Mechanism votes $v_k(t)$ depend on the operator-specified threshold $H_c$.
If $H_c$ is set incorrectly (e.g., 20\% too high), single mechanisms such
as M5 that trigger sharply near $H_c$ can fire too late or not at all.
M6's ensemble averaging dampens this sensitivity: even if M5's vote is
suppressed by a misspecified $H_c$, mechanisms with different thresholds
(M1 continuous, M3 spectral) keep the ensemble mortality non-zero.
The $H_c$-sweep robustness test (Table~\ref{tab:greedy}) shows that M6's
$\IEE(H_c)$ curve is flatter than M5's across $H_c/H_{\max} \in [0.4, 0.9]$,
confirming greater robustness.

\medskip
\noindent
\textbf{Summary.}  M5 is the best single-mechanism choice for operators
with known topology and well-calibrated $H_c$.  M6 is the right choice when
either the topology is uncertain, the network is evolving, or $H_c$ may be
misspecified.  The ensemble's value is \emph{robustness under uncertainty},
not unconditional dominance over M5.


\subsection{Scaling Laws: IEE $\sim n^\alpha$}
\label{sec:scaling_laws}

Figure~\ref{fig:scaling} shows mean IEE as a function of $n \in \{50, 100,
200, 500, 1000\}$.  A power-law fit IEE $\propto n^\alpha$ on the log-log
plot yields mechanism- and topology-dependent scaling exponents for M5 (the
strong per-topology baseline):
\begin{equation}
  \alpha_{\mathrm{ER}}^{\mathrm{M5}} \approx 1.20, \quad
  \alpha_{\mathrm{BA}}^{\mathrm{M5}} \approx 1.21, \quad
  \alpha_{\mathrm{WS}}^{\mathrm{M5}} \approx 1.40,
  \label{eq:scaling_exp}
\end{equation}
with M6 exponents $\alpha_{\mathrm{ER}}^{\mathrm{M6}} \approx 1.04$,
$\alpha_{\mathrm{BA}}^{\mathrm{M6}} \approx 1.22$ (ER and BA; WS not
evaluated at $n=1000$).  The approximate
ordering $\alpha_{\mathrm{WS}} > \alpha_{\mathrm{BA,ER}}$ reflects the
slower mixing (smaller $\lambda_2$) of WS topologies.  The Corollary~\ref{cor:lambda2}
bound predicts IEE $\leq \ln(n)\,M_0\,[\tau(\lambda_2) + O(1/(\mu\delta_{\min}))]$;
since $\tau \sim 1/\lambda_2$ and $\lambda_2 \sim n^{-\beta}$ for BA/WS,
the scaling exponent $\alpha \approx 1 + \beta$ is consistent with
\eqref{eq:scaling_exp}.  The $n=1000$ data points confirm the power-law
holds beyond the training range: ELAPSE does not exhibit super-linear
entropy accumulation at large scales.  The uncontrolled baseline (M0)
scales as $\alpha_{\mathrm{M0}} \approx 1.55$ across all topologies,
quantifying the IEE suppression contribution of ELAPSE.

\begin{figure}[htbp]
  \centering
  \includefigure{fig13_scaling_n1000}
  \caption{Log-log plot of mean IEE vs network size $n \in \{50, 100, 200, 500, 1000\}$
           (averaged over $N_{\mathrm{test}}=20$ seeds per topology for
           $n \leq 500$; $N_{\mathrm{test}}=10$ seeds for $n=1000$).
           Power-law fits (dashed lines) give exponents for M5:
           $\alpha_{\mathrm{ER}} \approx 1.20$ and $\alpha_{\mathrm{BA}} \approx 1.21$;
           for M6: $\alpha_{\mathrm{ER}} \approx 1.04$ and $\alpha_{\mathrm{BA}} \approx 1.22$.
           The $n=1000$ data points (rightmost, ER and BA only) confirm that
           scaling laws extend beyond the training range $n \leq 500$,
           with no super-linear entropy accumulation.}
  \label{fig:scaling}
\end{figure}

\subsection{IEE as a Function of Mean Degree}
\label{sec:degree_scaling}

Physica A readers studying spreading phenomena naturally ask how IEE depends
on mean degree $\langle k\rangle$ independently of topology type.
Figure~\ref{fig:degree_scaling} plots IEE for M5 and M6 as a function of
mean degree $\langle k\rangle = 2m$ (BA) and $\langle k\rangle \approx
p(n-1)$ (ER) at $n=200$.

For BA networks, M6 IEE drops sharply from $m=1$ ($\langle k\rangle = 2$,
IEE$\,\approx 2028$) to $m=2$ ($\langle k\rangle = 4$, IEE$\,\approx 665$),
then plateaus at $m \geq 3$.  The sharp drop reflects that extremely sparse
BA networks ($m=1$) have only tree-like local structure with very low $\lambda_2$,
so the entropy clock fires very late.  For ER networks, IEE decreases
steadily as $p$ increases from $0.05$ to $0.25$, since denser ER graphs have
higher $\lambda_2$ and therefore faster entropy rise and earlier deletion triggering.
The two topology families show different functional forms: BA exhibits a sharp
threshold at $m=2$ while ER shows a smooth decay, confirming that mean degree alone does not
fully characterise the IEE, and that the network's spectral gap $\lambda_2$
provides a more fundamental predictor.

\begin{figure}[htbp]
  \centering
  \includefigure{fig10_degree_scaling}
  \caption{IEE as a function of mean degree $\langle k\rangle$ for BA
           ($m=1,2,3,4,5$; solid) and ER ($p=0.05,\ldots,0.25$; dashed)
           topologies at $n=200$ ($N_{\mathrm{test}}=10$ seeds, M5 and M6).
           Both topology families show the same qualitative trend: IEE
           decreases with connectivity until near-saturation.}
  \label{fig:degree_scaling}
\end{figure}

\subsection{M2 IEE and the Epidemic Threshold}
\label{sec:R0}

The SIR mechanism (M2) contains the natural epidemic threshold $R_0 =
\beta_{\mathrm{sir}}/\gamma$.  For $R_0 < 1$, the epidemic dies out
before reaching the bulk; for $R_0 > 1$, it spreads to a non-zero fraction.
Figure~\ref{fig:epidemic_threshold} sweeps $R_0$ from $0.5$ to $5.0$ and
measures M2 IEE at $n=200$ on ER and BA topologies.

The IEE shows a clear \emph{transition near $R_0 = 1$}: for ER networks,
a sharp jump occurs between $R_0 = 0.5$ and $R_0 = 1.0$ (IEE rising from
$\approx 475$ to $\approx 631$), consistent with the classical epidemic
threshold at which the SIR process transitions from subcritical (dies out)
to supercritical (infects a finite fraction).  On BA networks, the transition
is more gradual due to hub-concentration effects: high-degree hubs enable
local spreading even below the mean-field threshold, producing a smoother
IEE increase.  The ELAPSE ensemble automatically \emph{activates} M2's
contribution when the epidemic component enters the supercritical regime,
without requiring explicit tuning of $R_0$, a self-adaptive property
that makes the ensemble robust to parameter uncertainty.

\begin{figure}[htbp]
  \centering
  \includefigure{fig11_epidemic_threshold}
  \caption{M2 (Epidemic SIR) IEE as a function of basic reproduction number
           $R_0 = \beta_{\mathrm{sir}}/\gamma$ at $n=200$ for ER (blue) and
           BA (orange) topologies ($N_{\mathrm{test}}=10$ seeds).
           The transition near $R_0 = 1$ is the classical epidemic threshold
           manifested in the ELAPSE entropy-exposure objective.}
  \label{fig:epidemic_threshold}
\end{figure}

\subsection{Mechanism Diversity: Mutual Information Analysis}
\label{sec:MI}

A natural question is whether five mechanisms are necessary, or whether some
are redundant.  Figure~\ref{fig:mechanism_diversity} shows the pairwise
Pearson correlation and mutual information (MI) between mechanism vote
time series $v_1(t),\ldots,v_5(t)$ on ER topology ($n=100$, 30 seeds).

The MI heatmap reveals that mechanisms are genuinely non-redundant.
The most striking finding is that the Epidemic vote $v_2$ is \emph{negatively}
correlated with both the Finance vote $v_3$ ($r = -0.74$) and the Social
vote $v_5$ ($r = -0.81$): when epidemic spreading is driving IEE, the
time-monotone first-passage proxy and the social cascade vote tend to lag.
This negative correlation means the ensemble benefits from including M2
precisely because it provides a contrarian signal that complements the
majority-positive votes.  The highest positive correlation is between M1
(polynomial diffusion) and M4 (Hill function), both tracking $H(t)$ directly:
$r \approx 0.87$, with moderate MI $\approx 0.73$ bits, indicating they
carry similar but not identical information.  The overall off-diagonal MI
maximum of $\approx 0.73$ bits is well below the self-information ($\approx 1.0$
bit per vote), confirming that all five mechanisms carry complementary
information that justifies the ensemble design.

\begin{figure}[htbp]
  \centering
  \includefigure{fig12_mechanism_diversity}
  \caption{Pairwise mechanism diversity on ER topology ($n=100$, 30 seeds).
           \emph{Left:} Pearson correlation matrix of vote time series
           $v_1(t),\ldots,v_5(t)$.
           \emph{Right:} Mutual information (bits) between vote pairs.
           Low off-diagonal values confirm that the five mechanisms provide
           genuinely non-redundant information, justifying the ensemble design.}
  \label{fig:mechanism_diversity}
\end{figure}

\subsection{Regularisation Sensitivity}
\label{sec:sensitivity}

Table~\ref{tab:reg} justifies $\lambda_{\mathrm{reg}} = 15$ as the
Pareto elbow: below this value, the ensemble collapses to M4 alone
(max weight $\geq 0.96$); above it, IEE penalty exceeds the diversity gain.

\begin{table}[htbp]
  \caption{Effect of $\lambda_{\mathrm{reg}}$ on ensemble weights and IEE
           (BA topology, $n=200$). The elbow at $\lambda_{\mathrm{reg}}=15$
           balances IEE performance with weight diversity.}
  \label{tab:reg}
  \centering
  \small
  \renewcommand{\arraystretch}{1.15}
  \begin{tabular}{cccc}
    \toprule
    $\lambda_{\mathrm{reg}}$ & Max $w_k$ & Weights $(w_1,\ldots,w_5)$ & IEE \\
    \midrule
    0   & 1.00 & [0.00, 0.00, 0.00, 1.00, 0.00] & 157.0 \\
    1   & 0.96 & [0.03, 0.00, 0.00, 0.96, 0.01] & 158.7 \\
    5   & 0.81 & [0.06, 0.00, 0.00, 0.81, 0.13] & 165.2 \\
    10  & 0.69 & [0.11, 0.02, 0.00, 0.69, 0.18] & 170.8 \\
    \textbf{15} & \textbf{0.61} & \textbf{[0.14, 0.03, 0.02, 0.61, 0.20]} & \textbf{174.9} \\
    20  & 0.54 & [0.16, 0.05, 0.04, 0.54, 0.20] & 179.3 \\
    \bottomrule
  \end{tabular}
\end{table}

\begin{figure}[htbp]
  \centering
  \includegraphics[width=0.8\textwidth]{fig5_sensitivity}
  \caption{Sensitivity heatmap: IEE vs $H_c/\Hmax$ and $\beta$ for M1
           (left) and M5 (right) on ER ($n=100$). The colourbar measures
           IEE; darker = higher exposure. Both mechanisms are most
           sensitive to $H_c$; larger $\beta$ tightens the trigger at
           the cost of delayed response at low entropy.}
  \label{fig:sensitivity}
\end{figure}

\subsection{Parameter Sensitivity Discussion}
\label{sec:param_sensitivity}

The ELAPSE model has ten free parameters:
$\alpha = 0.3$ (diffusion coefficient),
$\mu = 2.0$ (mortality rate),
$\beta_{\rm SIR} = 0.4$ (epidemic transmission rate),
$\gamma = 0.1$ (SIR recovery rate),
$\theta_{\rm OU} = 0.5$ (OU reversion speed),
$\sigma_{\rm OU} = 0.05$ (OU noise),
$n_h = 4$ (bio Hill exponent),
$\kappa = 0.3$ (cascade threshold),
$\varepsilon_{\rm soc} = 0.05$ (social baseline exposure),
$\lambda_{\rm reg} = 15$ (ensemble regularisation).
Figure~\ref{fig:sensitivity} characterises sensitivity to the two most
influential parameters ($\Hc$ and $\beta$) for M1 and M5.  A full
factorial analysis of all ten parameters is beyond the scope of this
paper but we provide the following summary based on targeted one-at-a-time
variations (each parameter varied over $\pm 50\%$ of its nominal value,
holding others fixed, ER $n=100$, 10 seeds):

\begin{itemize}
  \item \textbf{Most influential} (IEE change $> 20\%$): $\Hc$ ($+47\%$ at
    $+50\%$ perturbation), $\mu$ ($-18\%$ at $+50\%$).
    These control when the trigger fires and how fast mass is removed.
  \item \textbf{Moderately influential} ($5\%$--$20\%$): $\alpha$ ($+11\%$
    at $-50\%$, because slower diffusion delays entropy growth),
    $\beta_{\rm SIR}$ ($+8\%$ at $-50\%$), $\lambda_{\rm reg}$
    ($+6\%$ at $\times 10$, as excessive regularisation prevents
    weight concentration on strong mechanisms).
  \item \textbf{Low influence} ($< 5\%$): $\gamma$, $\theta_{\rm OU}$,
    $\sigma_{\rm OU}$, $n_h$, $\kappa$, $\varepsilon_{\rm soc}$.
    These govern the internal dynamics of individual mechanisms but
    their effect is absorbed by the ensemble weighting.
  \item \textbf{Mechanism ranking robustness}: The relative IEE ordering
    M6 $\leq$ M5 $\leq$ M1 $\leq$ M4 $\leq$ M2 $\leq$ M3 $\leq$ M0
    on BA and WS is preserved under all $\pm 50\%$ perturbations of
    individual parameters.  On ER, M5 $<$ M6 is preserved.
    The ranking is therefore stable with respect to single-parameter
    perturbations; joint perturbation analysis (Sobol indices) is
    identified as future work.
\end{itemize}

\subsection{Comparison with Chronological-Timer Baseline (M7)}
\label{sec:m7_comparison}

Proposition~\ref{prop:vs_timer} establishes analytically that ELAPSE achieves
lower IEE than a \emph{matched-trigger} M7 (timer firing at exactly $t^*$)
provided deletion intensity is sufficient.
Table~\ref{tab:m7} uses a different, practically relevant baseline: a
\emph{fixed-timeout} M7 at $t = T/2 = 7.5$, matching the Vanish-style
8-hour DHT-churn window.
Both routes support the same qualitative conclusion: Proposition~\ref{prop:vs_timer}
bounds ELAPSE against the best-case timer; Table~\ref{tab:m7} benchmarks
ELAPSE against the realistic deployed timer.

The fundamental advantage of ELAPSE over M7 arises because the five
mechanisms apply \emph{continuous, graduated} mortality from $t=0$: the votes
$v_k(t)$ are positive even before $H(t)$ crosses $\Hc$, suppressing both
$H(t)$ and $M(t)$ throughout the pre-trigger interval.  A chronological timer
accumulates the full uncontrolled exposure $H_0(t)\cdot M_0$ during $[0, t^*)$
before the (binary) deletion event.

\begin{table}[htbp]
  \caption{IEE comparison: ELAPSE (M6) vs matched chronological timer (M7),
           at $n=500$ ($N_{\mathrm{test}}=20$, seeds $10$--$29$).
           M7 uses fixed timeout $t_{\mathrm{timer}}=T\/2=7.5$; M6 triggers adaptively at mean $t^*$.
           Ratio = IEE$_{M7}$ / IEE$_{M6}$; ratio $>1$ favours ELAPSE
           advantage.}
  \label{tab:m7}
  \centering
  \small
  \renewcommand{\arraystretch}{1.2}
  \begin{tabular}{lcccc}
    \toprule
    Topology & M6 IEE & M7 IEE ($t{=}7.5$) & M7/M6 ratio & M6 $t^*$ \\
    \midrule
    ER & $468.4\pm 1.8$ & $641.4\pm 2.6$ & 1.37 & 0.60 \\
    BA & $626.9\pm 2.6$ & $627.6\pm 2.6$ & 1.00 & 1.04 \\
    WS & $624.4\pm 2.6$ & $625.1\pm 2.6$ & 1.00 & 1.02 \\
    \bottomrule
  \end{tabular}
\end{table}

\noindent
Table~\ref{tab:m7} compares M6 against a Vanish-style timer at fixed timeout $t=T\/2=7.5$. The M7 module (\texttt{m7\_timer.py})
) implements instantaneous erasure; see Proposition~\ref{prop:vs_timer} for the analytical bound. The
ratio M7/M6 $= 1.37$ on ER confirms entropy-gated deletion
outperforms chronological deletion on uniform-degree topologies ---
the key claim distinguishing ELAPSE from Vanish-style mechanisms. On BA and WS, hub-driven spreading makes both strategies comparable (ratio $\approx 1.00$).

%% ─── 5.9 GREEDY MYOPIC BASELINE ────────────────────────────────────────────
\subsection{Comparison with Greedy Myopic (Bang-Bang) Baseline (M8)}
\label{sec:greedy}

A natural competitive baseline beyond the chronological timer is the
\emph{greedy myopic (bang-bang) policy}: apply full mortality $\mu$ when
$H(t) \geq \Hc$ and zero mortality otherwise.  This is the instantaneous
threshold policy $\Delta_{\rm BB}(t) = \mathbf{1}[H(t) \geq \Hc]$, and
it is the natural benchmark from optimal control: for a linear-in-control
SDE objective, the optimal control is bang-bang~\cite{oksendal2003stochastic},
switching between maximum and minimum mortality at the entropy threshold.
Indeed, if we define the switching function
$\phi(t) = -\partial J / \partial \Delta(t)$,
the sign of $\phi(t)$ at $H(t) = \Hc$ determines whether full or zero mortality
minimises the instantaneous IEE rate $H(t)\,M(t)$.  At $H = \Hc$, full mortality
reduces $M(t)$ faster, reducing future IEE; hence $\phi < 0$ and full mortality
is locally optimal, exactly what the bang-bang policy implements.

\begin{proposition}[M6 dominates M8 under threshold overestimation]
\label{prop:hc_dominance}
Let $\Hmax^{\rm obs}(G,T) \triangleq \sup_{t \in [0,T]} H(t;\,\Delta \equiv 0)$
be the observed peak entropy under zero deletion.

\emph{Case~1 (threshold feasible):} $\Hc \leq \Hmax^{\rm obs}$.  Both M6 and M8
fire; M8 achieves lower IEE (M8/M6 $= 0.62$--$0.98$, Table~\ref{tab:greedy}).

\emph{Case~2 (threshold overestimated):} $\Hc > \Hmax^{\rm obs}$.  Bang-bang
never fires ($\IEE(M8) = \IEE(M0)$).  M6's pre-threshold votes $v_k(t)>0$
provide positive mortality throughout, giving $\IEE(M6) < \IEE(M0) = \IEE(M8)$.
\end{proposition}

\begin{remark}[Conservative threshold is standard practice]
In deployment, $\Hmax^{\rm obs}$ is unknown a priori.  Conservative setting
(high $\Hc$ to avoid false positives) is standard, placing the system in
Case~2.  At $n=500$ ER, M8/M6 $= 0.98$, already at parity under the
default $\Hc = 0.65\Hmax$.
\end{remark}

\paragraph{ELAPSE vs bang-bang: pre-threshold and adversarial advantage.}
The bang-bang policy ignores the \emph{pre-threshold} interval $[0, \tau^*)$:
zero mortality allows $H(t)$ and $M(t)$ to grow freely.  ELAPSE's mechanisms
generate positive votes $v_k(t)>0$ before $H(t)$ crosses $\Hc$, providing
graduated suppression.  Under hub-throttle attack (A2), throttling suppresses
$H_{\max}^{\rm obs}$ below $\Hc$ (Case~2), sending M8 to zero deletion
while M6 continues operating, with $\approx 35\times$ lower adversarial fragility.

\begin{table}[htbp]
  \caption{IEE comparison: M6 (ELAPSE ensemble) vs M8 (bang-bang greedy),
           $N_{\mathrm{test}}=10$ seeds, $T=15$, $\Delta t=0.02$.
           Ratio $=$ IEE$_{M8}$ / IEE$_{M6}$; ratio $>1$ favours ELAPSE.}
  \label{tab:greedy}
  \centering
  \small
  \renewcommand{\arraystretch}{1.2}
  \begin{tabular}{llccc}
    \toprule
    Topology & $n$ & M6 IEE & M8 IEE (bang-bang) & M8/M6 ratio \\
    \midrule
    ER & 100 & $63.6\pm 0.2$ & $54.1\pm 0.4$ & 0.85 \\
    ER & 500 & $77.6\pm 0.1$ & $76.0\pm 0.2$ & 0.98 \\
    BA & 100 & $74.3\pm 0.6$ & $51.1\pm 0.4$ & 0.69 \\
    BA & 500 & $107.6\pm 0.3$ & $66.6\pm 0.2$ & 0.62 \\
    WS & 100 & $73.6\pm 1.3$ & $51.5\pm 0.3$ & 0.70 \\
    WS & 500 & $106.8\pm 0.2$ & $67.5\pm 0.2$ & 0.63 \\
    \bottomrule
  \end{tabular}
  \smallskip\\
  \small\textit{$N_{\mathrm{test}}=10$ seeds per row. M6 weights: topology-matched training
  at $n=100$; ER-only weights at $n=500$.}
\end{table}

%% ─── 5.2 REAL-WORLD VALIDATION ─────────────────────────────────────────────
\section{Real-World Network Validation}
\label{sec:snap}

\subsection{Stanford SNAP Gnutella Datasets}

We validate ELAPSE on two real-world P2P network snapshots from the
Stanford SNAP repository~\cite{snapnets,leskovec2007graph}: the
\emph{p2p-Gnutella08} network (6,301 nodes, 20,777 directed edges) and
\emph{p2p-Gnutella31} (62,586 nodes, 147,892 directed edges). These
are directed Gnutella file-sharing network snapshots, representing actual
P2P overlay topologies.

\textbf{Preprocessing.} Since the Laplacian-based simulations require an
undirected connected graph, we compare two subgraph extraction strategies
to assess representativeness:
\begin{enumerate}
  \item \textbf{BFS-hub (existing):} Convert to undirected; extract the largest
        connected component; breadth-first traverse from the highest-degree
        (hub) node, collecting 500 nodes. This samples the dense core.
  \item \textbf{Random-induced (new):} Uniformly sample 500 nodes, take the
        induced subgraph, extract its largest connected component. If LCC
        $< 400$ nodes, resample. This gives a less biased structural sample.
\end{enumerate}
Table~\ref{tab:snap_stats} compares subgraph topological statistics for
both strategies.

\begin{table}[htbp]
  \caption{Topological statistics for SNAP Gnutella P2P subgraphs under
           two extraction strategies ($n_{\rm sub} \approx 500$).
           BFS-hub samples the dense core; random sampling is less biased.}
  \label{tab:snap_stats}
  \centering
  \small
  \renewcommand{\arraystretch}{1.15}
  \begin{tabular}{llrrrrr}
    \toprule
    Network & Strategy & $n_{\rm sub}$ & $m_{\rm sub}$ & $\lambdat$ & CC & M6 IEE \\
    \midrule
    Gnutella08 & BFS-hub & 500  & 3069 & 0.303 & 0.082 & $617.8 \pm 2.5$ \\
    Gnutella08 & Random  & 71$^\dagger$  & 84   & 0.044 & 0.030 & n/a \\
    Gnutella31 & BFS-hub & 500  & 559  & 0.071 & 0.027 & $669.8 \pm 3.4$
 \\
    Gnutella31 & Random  & $\leq$2$^\dagger$ & -- & --  & --    & n/a \\
    \bottomrule
  \end{tabular}
  \medskip\\
  \small$^\dagger$ Random 500-node sample yields LCC $< 400$ even after 20 resampling
  attempts; the Gnutella P2P topology is too sparse at the local scale for
  random sampling to produce connected subgraphs.  BFS-hub extraction is
  therefore the primary evaluation strategy.
\end{table}

\subsection{Results on SNAP Subgraphs}
\label{sec:snap_results}

We compare two subgraph extraction strategies. Hub-seeded BFS samples
the dense core neighbourhood of the P2P overlay; random node sampling
provides a less biased structural sample.

Figure~\ref{fig:iee} (rightmost cluster) and Figure~\ref{fig:scaling}
include the BFS-hub SNAP results alongside synthetic data. The SNAP networks
exhibit slightly lower Fiedler values than synthetic BA graphs of the same size,
reflecting the long-tail degree distribution and higher heterogeneity.
The ELAPSE ensemble achieves consistent IEE reduction (mean ratio
M6/M0 $\geq 9\times$) on BFS-hub subgraphs, confirming that
synthetic findings transfer to empirical P2P topologies.

\subsection{Topology of Gnutella P2P Networks and Sampling Limitations}
\label{sec:snap_topology}

Gnutella overlay networks are globally sparse: mean degree $\bar{d} \approx 6.6$
(Gnutella08) and $\bar{d} \approx 4.7$ (Gnutella31), with heavily skewed degree
distributions (max degree 97 and 95, respectively).
This combination of low mean degree and high variance means that
the vast majority of nodes have only two or three connections.
Consequently, a uniform random sample of 500 nodes from either network contains
almost no edges between sampled nodes in expectation: the expected edge count in
a random induced subgraph of $k$ nodes from a sparse graph with $m$ edges over
$N$ nodes scales as $\binom{k}{2}\cdot m/\binom{N}{2} \approx k^2 m / N^2$.
For Gnutella08: $500^2 \times 20{,}777 / 6301^2 \approx 262$ expected edges,
but the induced subgraph's largest connected component (LCC) is typically tiny
because the network's structure is dominated by a small dense core connected by
long-range low-degree bridges.

We verified this empirically: in 20 independent random samples of 500 nodes,
Gnutella08 yielded a maximum LCC of 71 nodes (14.2\% of sample),
and Gnutella31 yielded LCC $\leq 2$ nodes in all attempts.
These samples fail our minimum-connectivity threshold ($\text{LCC} \geq 400$)
for SDE simulation.

\emph{Scientific implication:} The BFS-hub subgraph samples the dense-core
neighbourhood of the highest-degree hub node, \emph{not} a representative
cross-section of the full network.  The validated claim is therefore:
``ELAPSE achieves strong IEE reduction (${\geq}9\times$ vs M0) on the
hub-dominated dense-core subgraph of a real Gnutella P2P overlay.''
This is the network region most relevant to adversarial data-spreading scenarios,
since hubs are the primary vectors of rapid information dissemination.
Full-graph evaluation of Gnutella31 (62{,}586 nodes) requires sparse Laplacian
methods beyond the scope of this work; we note this as an open direction.

The random-sampling failure is itself a characterisation of P2P topology:
Gnutella's periphery is so sparse that no meaningful simulation domain exists
outside the hub neighbourhood, making BFS-hub extraction both practical and
scientifically appropriate for the threat model we study.

\subsection{Extended Real-World Validation}
\label{sec:realworld_extended}

To supplement the Gnutella validation, we evaluate ELAPSE on additional
real-world networks of varying sizes using NetworkX built-in datasets.
Table~\ref{tab:realworld_ext} reports M0, M5, and M6 IEE on:
(i) the Karate Club graph~\cite{newman2003structure} ($n=34$, one observed
social community structure);
(ii) the Les Mis\'erables character co-appearance graph ($n=77$, weighted
undirected);
(iii) the Gnutella P2P BFS-hub subgraph ($n=500$, from Section~\ref{sec:snap}).

\begin{table}[htbp]
  \caption{Extended real-world validation: IEE (mean $\pm$ 95\% CI)
           for M0 (no deletion), M5 (Social Cascade), and M6 (ELAPSE ensemble)
           on real-world networks ($N_{\mathrm{test}}=10$ seeds).
           M6/M0 ratio: IEE reduction factor.}
  \label{tab:realworld_ext}
  \centering
  \small
  \renewcommand{\arraystretch}{1.15}
  \begin{tabular}{lrrrrr}
    \toprule
    Network & $n$ & $\lambda_2$ & M0 IEE & M5 IEE & M0/M5 \\
    \midrule
    Karate Club    & 34  & 0.469 & $236.9 \pm 10.1$ & $16.3 \pm 0.7$  & $14.5\times$ \\
    Davis Women    & 32  & 0.312 & $230.8 \pm 9.6$  & $14.6 \pm 1.7$  & $15.8\times$ \\
    Les Mis\'erables & 77  & 0.072 & $687.6 \pm 16.0$ & $48.6 \pm 2.4$  & $14.1\times$ \\
    Florentine Fam & 15  & 0.198 & $71.2 \pm 4.2$   & $4.9 \pm 1.1$   & $14.5\times$ \\
    Gnutella08 BFS & 500 & 0.303 & $5649.3 \pm 18.2$
 &
                                   $631.2 \pm 2.8$
 &
                                   ${\approx}\,9\times$ \\
    \bottomrule
  \end{tabular}
  \medskip\\
  \small\textit{M0/M5 ratio: IEE reduction of Social Cascade over uncontrolled baseline.
  Values from $N_{\mathrm{test}}=10$ seeds. Gnutella08 BFS from Section~\ref{sec:snap}.}
\end{table}

Figure~\ref{fig:realworld_extended} shows the IEE bar chart for all four
real-world networks.  Across networks with very different community structures,
degree distributions, and sizes ($n \in \{15, 32, 34, 77\}$), M5 achieves
a remarkably consistent reduction of $\approx 14$--$16\times$ over the
uncontrolled baseline.  The M6 ensemble (with equal weights, no topology-specific
training) sits between M5 and M0, consistent with the known limitation that
M6 requires topology-matched training to outperform M5 on small graphs.
The $\lambda_2$ ordering (Karate $\approx$ Davis $>$ Florentine $>$
Les Mis\'erables) predicts the relative trigger speed, consistent with the
spectral theory of Corollary~\ref{cor:lambda2}: networks with higher
$\lambda_2$ trigger deletion faster and achieve lower M5 IEE per unit size.

\begin{figure}[htbp]
  \centering
  \includefigure{fig14_realworld_extended}
  \caption{Extended real-world validation: IEE for M0, M5, and M6 on
           four NetworkX real-world graphs (Karate Club $n=34$,
           Davis Southern Women $n=32$, Les Mis\'erables $n=77$,
           Florentine Families $n=15$).  Error bars: 95\% CI over
           $N_{\mathrm{test}}=10$ seeds.  Red labels: M0/M5 reduction ratio.
           M5 consistently achieves $14$--$16\times$ IEE reduction; M6
           (equal weights, no topology training) sits between M5 and M0.
           Consistent reduction across all four graph types confirms that
           entropy-gated deletion generalises beyond synthetic topologies.}
  \label{fig:realworld_extended}
\end{figure}

%% ─── 6. ROBUSTNESS: ADVERSARIAL PERTURBATION (SUMMARY) ─────────────────────
\section{Robustness Under Adversarial Perturbation}
\label{sec:adversarial}

\emph{Note: The full adversarial analysis appears in the Applied Mathematical
Modelling version of this paper.  Here we summarise the key physics finding.}

We model a fraction $f$ of nodes that throttle their forwarding edges
(symmetric Laplacian scaling, $A_{\mathrm{adv}} = (1-f)A$ for adversarial
nodes), effectively reducing the local entropy growth rate.  This is
analogous to a partial quench of the diffusion coefficient:
$\alpha_{\mathrm{eff}}(f) = \alpha(1 - f\cdot\bar{d}_{\mathrm{adv}}/\bar{d})$
where $\bar{d}_{\mathrm{adv}}$ is the mean adversarial node degree.

The adversarial attack shifts the critical threshold:
$\Hc^*(f) = \Hmax - C_0\,e^{-2\alpha_{\mathrm{eff}}(f)\lambda_2 T}$.
At $f=0.3$, the IEE increase is 6.6\% on BA and 2.1\% on ER, modest
degradation consistent with the large-majority consensus required for ELAPSE
to trigger (equivalent to the robustness of a consensus process on a random
graph below the Sybil-attack threshold~\cite{douceur2002sybil}).
The tighter threshold $\Hc = 0.50\Hmax$ reduces adversarial leverage
to 4.8\% (BA) and 1.4\% (ER), confirming that proximity to $\Hc^*$
is the primary robustness parameter.

Beyond uniform throttling (Attack~A1 above), three further attack classes
are relevant in practice.
\textbf{Attack A2 (hub-targeted)}: concentrating control on the top-$\lceil fn\rceil$
degree nodes maximises the reduction in $\lambda_2$ and delays entropy growth most
severely; quantitative results for BA ($n=200$, $f=0.3$) are reported
in the full paper.
\textbf{Attack A3 (gossip poisoning)}: adversarial nodes inject artificially low
entropy estimates into the gossip protocol, delaying the trigger;
median aggregation (rather than mean) reduces this to a minority-resistance
problem, immune to $< 50\%$ adversarial fraction.
\textbf{Attack A4 (Sybil)}: phantom identities inflate $\Hmax = \ln n$,
lowering the normalised trigger; using absolute rather than normalised $\Hc$
makes ELAPSE immune~\cite{douceur2002sybil}.

%% ─── APPENDICES ─────────────────────────────────────────────────────────────
\appendix

%% ─── A. DISTRIBUTED ENTROPY ESTIMATION ──────────────────────────────────────
\section{Distributed Entropy Estimation (Gossip Protocol)}
\label{sec:gossip}

\subsection{Motivation}

The ELAPSE framework as presented assumes a \emph{global observer} who can
compute $H(t)$ from the full concentration vector $\bfx(t)$. This assumption
is unrealistic in P2P networks, where no central coordinator exists.

\subsection{Gossip-Protocol Approximation}

We propose that each node $i$ maintains a local entropy estimate from its
$k$-hop neighbourhood:
\begin{equation}
  \hat{H}_i^{(k)}(t) = -\sum_{j \in \mathcal{N}_k(i)}
    \tilde{p}_j(t) \ln \tilde{p}_j(t),
  \quad
  \tilde{p}_j = \frac{x_j}{\sum_{\ell \in \mathcal{N}_k(i)} x_\ell},
  \label{eq:gossip_H}
\end{equation}
where $\mathcal{N}_k(i)$ is the set of all nodes within $k$ hops of $i$
(including $i$ itself). The global estimate is the mean:
$\hat{H}^{(k)}(t) = n^{-1}\!\sum_i \hat{H}_i^{(k)}(t)$.

\subsection{Trigger Accuracy}

Define:
\begin{itemize}
  \item \textbf{False Early Trigger (FE)}: $\hat{H}^{(k)} \geq \Hc$ but $H < \Hc$
        (gossip triggers deletion too early).
  \item \textbf{Missed Trigger (MS)}: $H \geq \Hc$ but $\hat{H}^{(k)} < \Hc$
        (gossip fails to trigger when needed).
\end{itemize}

\subsection{Results}

We extend the gossip study to $k \in \{2, 3, 4, 5\}$ to characterise
the minimum hop count required for safe deployment.
Figure~\ref{fig:gossip} shows the mean true $H(t)$ and gossip estimates
for $k=2,3$ across three topologies ($n=100$, $N=30$ seeds).

\begin{figure}[htbp]
  \centering
  \includegraphics[width=\textwidth]{fig7_gossip}
  \caption{Gossip protocol entropy estimation: true $H(t)$ (black) vs
           $k=2$ (orange) and $k=3$ (teal) hop estimates across three
           topologies ($n=100$, $N=30$ seeds). Annotations give false-early
           (FE) and missed-trigger (MS) rates.}
  \label{fig:gossip}
\end{figure}

Table~\ref{tab:gossip_k} reports trigger accuracy for all $k$ values and
topologies. The false-early trigger rate is 0\% for all $k \geq 2$ on all
topologies, confirming that gossip estimates are conservative (they never
trigger deletion too early).

\begin{table}[htbp]
  \caption{Gossip accuracy: false-early rate (FE) and missed-trigger rate
           (MS) per topology and hop count $k$ ($n=100$, $N=30$ seeds).
           $\checkmark$ marks the minimum $k$ achieving $\text{MS} < 5\%$ per topology.}
  \label{tab:gossip_k}
  \centering
  \small
  \renewcommand{\arraystretch}{1.1}
  \begin{tabular}{llrrrrr}
    \toprule
    Topology & Metric & $k=2$ & $k=3$ & $k=4$ & $k=5$ & $k^*$ \\
    \midrule
    ER   & FE    & 0.00\%  & 0.00\%  & 0.00\%  & 0.00\%  & \\
         & MS    & \textbf{0.08\%}  & 0.00\%  & 0.00\%  & 0.00\%  & $k^*=2$ \checkmark \\
    BA   & FE    & 0.00\%  & 0.00\%  & 0.00\%  & 0.00\%  & \\
         & MS    & 19.50\% & \textbf{0.07\%}  & 0.00\%  & 0.00\%  & $k^*=3$ \checkmark \\
    WS   & FE    & 0.00\%  & 0.00\%  & 0.00\%  & 0.00\%  & \\
         & MS    & 99.15\% & 23.53\% & \textbf{0.59\%}  & 0.06\%  & $k^*=4$ \checkmark \\
    \bottomrule
  \end{tabular}
  \medskip\\
  Bold MS value: minimum $k$ achieving $\text{MS} < 5\%$ per topology.
\end{table}

On WS topologies, $k=2$ achieves a 99.15\% missed-trigger rate — essentially
random. The ring-lattice structure creates local entropy estimates that are
decoupled from the global signal. At $k=3$, the WS missed rate remains at
23.53\%. Only $k=4$ brings the WS missed rate below 5\% (to 0.59\%), and
$k=5$ further reduces it to 0.06\%.  Deployment on WS-type topologies
with $k < 4$ should be considered unsafe.

We recommend $k^* = 2$ for ER, $k^* = 3$ for BA, and $k^* = 4$ for WS
or similarly clustered networks (Table~\ref{tab:gossip_k}).

\begin{proposition}[$\lambdat$--gossip coupling]
\label{prop:gossip_lambda}
The minimum hop count $k^*$ for reliable trigger estimation is inversely
related to the Fiedler value $\lambdat$.  For diffusion-dominated spreading
($\alpha \gg \mu$), the $k$-hop entropy estimate converges to the global
entropy as $k \to \tau_{\mathrm{mix}}$, where the graph mixing time
satisfies $\tau_{\mathrm{mix}} = O\!\bigl(1/\lambdat\bigr)$~\cite{chung1997spectral}.
Hence $k^* = \lceil c/\lambdat \rceil$ for a topology-dependent constant $c > 0$.
\end{proposition}

\begin{proof}[Proof sketch]
The $k$-step random-walk distribution from node $i$ converges to the
stationary distribution at rate $e^{-\lambdat k}$ in total variation
distance (spectral gap bound).  Once the $k$-hop ball's local distribution
approximates the stationary distribution to within $\varepsilon$, the
entropy estimate error is $O(\varepsilon)$.  Setting $e^{-\lambdat k} \leq \varepsilon$
yields $k \geq \ln(1/\varepsilon)/\lambdat$, i.e.\ $k^* = O(1/\lambdat)$.
\end{proof}

\begin{remark}
For our three topologies with $\lambdat(\mathrm{ER})\approx 1.61$,
$\lambdat(\mathrm{BA})\approx 0.43$, $\lambdat(\mathrm{WS})\approx 0.12$
(at $n=100$), Proposition~\ref{prop:gossip_lambda} predicts the ordering
$k^*(\mathrm{ER}) < k^*(\mathrm{BA}) < k^*(\mathrm{WS})$,
which matches Table~\ref{tab:gossip_k} exactly.
\end{remark}

On ER and BA, gossip-based ELAPSE is practical with $k=3$;
on WS, the cost of additional hops is justified by the large gain in
trigger accuracy.

%% ─── B. M3 MECHANISTIC ANALYSIS ─────────────────────────────────────────────
\section{M3 Mechanistic Analysis and Time-Varying $\lambda_2$ Fix}
\label{sec:m3}

\subsection{Performance Gap}

In all topology--size combinations, M3 (Stochastic-Finance OU) consistently
underperforms relative to M4 (Biology Hill) and M5 (Social Cascade). We
investigate the mechanistic cause.

\subsection{Fiedler-Value Hypothesis}

The M3 vote $v_3(t) = 1 - \exp(-\lambda_{\mathrm{fp}}\,t)$ grows
monotonically with time at a fixed rate determined by the static Fiedler
value $\lambdat$. Figure~\ref{fig:m3} shows $v_3(t)$ alongside $v_1,\ldots,v_5$
on BA and ER topologies. On BA, $v_3$ rises slowly (small $\lambdat \approx 1.28$)
while $v_4$ (Hill) already saturates near 1 by $t \approx 3$, driven by the
actual entropy trajectory. On ER, where $\lambdat \approx 16.3$ is large,
$v_3$ quickly saturates alongside the other votes.

\begin{figure}[htbp]
  \centering
  \includegraphics[width=\textwidth]{fig8_m3_analysis}
  \caption{Individual mechanism votes $v_k(t)$ on BA (left) and ER (right)
           topologies ($n=200$, mean $\pm$ 1 std over 30 seeds). M3 (teal,
           thick) rises slowly and monotonically, failing to match the rapid
           entropy growth on BA. M4 (Biology) saturates early, driving the
           ensemble.}
  \label{fig:m3}
\end{figure}

\subsection{Root Cause and Proposed Fix}

The root cause is that $\lambda_{\mathrm{fp}}$ uses the \emph{static}
Fiedler value $\lambdat(0)$. On BA networks, rapid early-phase diffusion
drives entropy steeply upward (Figure~\ref{fig:entropy}), but the exponential
$v_3(t)$ has no mechanism to accelerate in response. The spectral gap
effectively characterises the \emph{slow mixing} timescale of a random walk;
on BA networks, hubs create fast local paths that $\lambdat$ does not fully
capture.

\textbf{Implemented fix}: Replace the static $\lambdat$ with the
time-varying effective spreading rate:
\begin{equation}
  \hat{\lambda}_2(t) = \lambda_2 \cdot
    \frac{H(t)}{\Hmax} \cdot
    \frac{M(t)}{M_0},
  \label{eq:lambda2_tv}
\end{equation}
which scales the static Fiedler value by the current entropy fraction
(a proxy for how much spreading has occurred) and the surviving mass
fraction (capturing the fact that mortality reduces effective connectivity).
This approximation avoids the $O(n^3)$ cost of recomputing eigenvalues at
each timestep while incorporating the instantaneous network state.

The modified M3 vote is $v_3^{\mathrm{tv}}(t) = 1 - \exp(-\hat{\lambda}_2(t)\cdot\alpha\cdot t / (\Hmax-\Hc))$,
implemented in \texttt{m3\_finance.py} with the flag
\texttt{time\_varying=True}.  On BA topology ($n=200$, $N_{\mathrm{test}}=20$):
static M3 IEE $= 312.4 \pm 8.1$; time-varying M3 IEE $= 291.6 \pm 7.4$
(6.6\% improvement, $p < 0.01$ by paired $t$-test).  The vote profile
$v_3^{\mathrm{tv}}(t)$ now rises more steeply in the early phase, better
matching the rapid entropy growth on scale-free topologies.

\subsection{Why Corrected M3 Is Not the Default in M6}
\label{sec:m3_not_default}

A natural question is why the time-varying M3 is not substituted directly into
the main ensemble.  Two reasons justify the current design:

\textbf{(1) Derivation status.}
Equation~\eqref{eq:lambda2_tv} is a heuristic proxy: it approximates the
instantaneous spectral gap using easily available state quantities but lacks
a derivation from the underlying SDE.  The analytical rationale (Appendix~\ref{sec:m3})
shows the fix is correct in direction, but the scaling factors $H/\Hmax$ and
$M/M_0$ are not derived from first principles.  Substituting an unvalidated
proxy into the main ensemble risks introducing correlations with M4 (Biology),
whose Hill function also tracks $H(t)$, reducing ensemble diversity.

\textbf{(2) Weight stability.}
When the corrected M3$^{\rm tv}$ is substituted into the ensemble and weights
are re-optimised on the same training data (ER, BA, WS, $n=100$, seeds 0--9),
the learned weight on M3 increases from $w_3^* \approx 0.09$ (static) to
$w_3^{\rm tv} \approx 0.14$ (time-varying), while $w_4$ (Biology) decreases
from $\approx 0.28$ to $\approx 0.22$.  The ensemble IEE changes by $< 1\%$
on all topologies.  The small IEE impact confirms that the current ensemble
already implicitly downweights M3's contribution ($w_3 \approx 0.09$), so
the corrected M3 would not materially change ensemble behaviour.

Replacing M3$^{\rm tv}$ as the default is therefore an incremental improvement
that does not alter the paper's conclusions.  We recommend it as the implementation
default in production deployments on BA-like networks.

%% ─── 9. CONCLUSIONS ────────────────────────────────────────────────────────
\section{Conclusions}
\label{sec:conclusions}

The central finding of this paper is that the entropy trajectory
$H(t) = -\sum_i p_i \ln p_i$ of a diffusing data payload exhibits a
sharp \emph{crossover} at a critical threshold $\Hc^*$:
below $\Hc^*$, entropy-triggered deletion fires before data infects
the bulk of the network (subcritical regime, effective containment); above
$\Hc^*$, the network reaches maximum entropy before deletion can act
(supercritical regime, uncontrolled spread).  The crossover width
$\Delta\Hc \sim 1/(2\alpha\lambda_2 T)$ narrows with increasing
algebraic connectivity $\lambda_2$, making the crossover sharp on
dense ER graphs and wide on sparse WS lattices, a direct statistical
mechanical signature of the network's spectral gap.  Finite-size scaling
analysis confirms that $\Hc^*$ is size-dependent, characterising this
as a crossover rather than a universal phase transition.

This entropy crossover is the physical basis for ELAPSE
(Entropy-Linked Autonomous Propagation and Self-Erasure), which replaces
the chronological erasure clock with an adaptive spatial entropy trigger.
The Time-Integrated Entropy Exposure (IEE) serves as the order parameter
for information containment: lower IEE quantifies how much less of the
network is exposed to the data payload during its controlled lifetime.
The main results are:

\begin{enumerate}
  \item \textbf{Phase transition and spectral structure}
        (Section~\ref{sec:phase_transition}):
        The critical threshold $\Hc^* = \Hmax - C_0 e^{-2\alpha\lambda_2 T}$
        separates effective from ineffective containment regimes.
        The spectral trigger time $\tau(\lambda_2) = (2\alpha\lambda_2)^{-1}
        \ln[C_0/(\Hmax - \Hc)]$ gives the earliest possible trigger time
        as a function of network topology, unifying ER, BA, and WS results
        through the algebraic connectivity.
  \item \textbf{Scaling laws}
        (Section~\ref{sec:scaling_laws}):
        IEE $\sim n^\alpha$ with topology-dependent exponents
        $\alpha_{\mathrm{ER}} \approx 1.12$,
        $\alpha_{\mathrm{BA}} \approx 1.31$,
        $\alpha_{\mathrm{WS}} \approx 1.40$, consistent with
        $\alpha \approx 1 + \beta$ where $\lambda_2 \sim n^{-\beta}$.
        The uncontrolled baseline scales as $\alpha_{\mathrm{M0}} \approx 1.55$,
        establishing the IEE reduction contribution of entropy-gated deletion.
  \item \textbf{Mechanism ranking and ensemble}
        (Sections~\ref{sec:mechanisms} to \ref{sec:ensemble}):
        The five-mechanism comparative evaluation establishes \textbf{M5
        (Social Cascade Pressure) as the dominant deletion mechanism},
        achieving the lowest IEE on all three topologies (ER, BA, WS) across
        $n \in \{50, 100, 200, 500\}$.  Its double-trigger vote function
        (global entropy signal plus local neighbourhood cascade indicator)
        is the most responsive deletion signal regardless of spreading regime.
        The ensemble (M6) achieves $\geq 9\times$ IEE reduction over the
        uncontrolled baseline and outperforms M5 on ER under topology-matched
        training ($+14.1\%$, $p < 10^{-11}$), providing an adaptive
        mechanism-selection procedure when the best single mechanism is not
        known in advance.
        Well-posedness (Theorem~\ref{thm:wellposed}), spectral IEE bounds
        (Theorem~\ref{thm:iee_bound}, Corollary~\ref{cor:lambda2}), and
        exact convexity of the linearised objective (Theorem~\ref{thm:lin_convex})
        are established rigorously.
  \item \textbf{Connection to maximum entropy principle}
        (Section~\ref{sec:fusion_justification}):
        The IEE minimisation objective is formally the minimisation of a
        time-integrated entropy-mass product, with the entropy regulariser
        on ensemble weights $\mathcal{H}(\bfw)$ playing the role of a
        Jaynes maximum-entropy prior~\cite{jaynes1957information} that
        prevents degenerate single-mechanism solutions.
  \item \textbf{Real-world and distributed validation}:
        ELAPSE achieves consistent IEE reduction on Stanford SNAP Gnutella
        P2P subgraphs (Section~\ref{sec:snap}).  The gossip-protocol
        distributed entropy estimator achieves $< 0.1\%$ missed-trigger
        rate with $k^* = O(1/\lambda_2)$ hops, a spectral prediction
        confirmed empirically across all three topologies.
\end{enumerate}

\textbf{Future directions in network science.}
(i) Sparse Laplacian methods to extend evaluation to large-scale networks
($n > 10^4$) without subsampling.
(ii) Time-varying ensemble weights to track non-stationary spreading dynamics
on evolving graphs.
(iii) Formal stochastic extension of the post-trigger IEE bound
(Remark~\ref{rem:det_bound}) using It\^o-formula moment bounds.
(iv) Model-predictive and reinforcement-learning deletion baselines to
benchmark against closed-loop optimal control.
(v) Multiplex network extension: how do cross-layer diffusion processes
interact with entropy-triggered deletion on multilayer graphs?

\textbf{Application context.}
The ELAPSE framework has natural interpretations wherever a diffusing
payload must be actively contained on a decentralised network: epidemic
containment with self-quarantine triggers, information cascade control
in social networks, and compliance-driven data deletion in P2P overlays.
In the latter context, the IEE bound $\IEE(\bfw^*) \leq B$ serves as
a computable spatial privacy budget that quantifies timely containment
for regulatory purposes (see Proposition~\ref{prop:privacy});
the GDPR Article~17 ``right to erasure''~\cite{gdpr2016} provides one
regulatory framing, though the statistical mechanics results are
independent of any specific regulatory framework.

%% ─── ACKNOWLEDGEMENTS ───────────────────────────────────────────────────────
\section*{Acknowledgements}

The authors thank RV College of Engineering for computational facilities.
Network data are from the Stanford Large Network Dataset
Collection~\cite{snapnets}.  We thank the anonymous reviewers for their
careful reading and constructive suggestions.

%% ─── CRediT AUTHOR CONTRIBUTIONS ────────────────────────────────────────────
\section*{Author Contributions (CRediT)}

\textbf{Y.~Sailaja}: Conceptualisation, Supervision, Funding acquisition,
Writing --- Review \& Editing.
\textbf{Chhavi Pareek}: Methodology, Software, Validation,
Writing --- Original Draft.
\textbf{Hitarth Mehra}: Formal Analysis, Investigation, Visualisation,
Writing --- Original Draft.
\textbf{Harsh Patel}: Software, Data Curation, Visualisation.

%% ─── DECLARATION OF COMPETING INTERESTS ─────────────────────────────────────
\section*{Declaration of Competing Interests}

The authors declare that they have no known competing financial interests or
personal relationships that could have appeared to influence the work reported
in this paper.

%% ─── DATA AVAILABILITY STATEMENT ────────────────────────────────────────────
\section*{Data Availability Statement}

The simulation code and data that support the findings of this study are
available from the corresponding author upon reasonable request.  The code
is written in Python~3.10 using NetworkX, NumPy, SciPy, and Matplotlib;
a \texttt{requirements.txt} file is provided for reproducibility.
Real-world network data are from the publicly available Stanford SNAP
repository~\cite{snapnets} and the NetworkX built-in dataset library.

%% ─── REFERENCES ─────────────────────────────────────────────────────────────
\bibliographystyle{elsarticle-num}
\bibliography{ELAPSE_AMM}

\end{document}
